% ---------------------------------------------------
% KI-VERZEICHNIS
% Eingebunden in main.tex mit % ---------------------------------------------------
% KI-VERZEICHNIS
% Eingebunden in main.tex mit % ---------------------------------------------------
% KI-VERZEICHNIS
% Eingebunden in main.tex mit % ---------------------------------------------------
% KI-VERZEICHNIS
% Eingebunden in main.tex mit \input{ki_verzeichnis}
% Benoetigt in der Praeambel: \usepackage{xltabular}
% Regeln: Projekt-Notiz KI-Verzeichnis_Regeln_2026-09-30
%   Spalten: Kennung, Modell, Datum, Prompt, Verwendung
%   Kennung: Claude2026a ... Claude2026z, Claude2026aa, Claude2026ab ...
%   Sortierung: nach Vorkommen, erst Gesamttext, dann Abschnitte
%               und Literatur, dann Code in der Reihenfolge der Entstehung
%               (Pipeline: Aufbereitung, Vorpruefung, Modelle, README,
%               Kommentare)
%   Prompts ohne Anfuehrungszeichen, ohne Kursivschrift
% ---------------------------------------------------
\newpage
\phantomsection
\section*{KI-Verzeichnis}
\addcontentsline{toc}{section}{KI-Verzeichnis}

\begingroup
\tabformat
\setlength{\tabcolsep}{4pt}
\begin{xltabular}{\textwidth}{@{}%
    >{\raggedright\arraybackslash\hyphenpenalty=10000}p{2.6cm}%
    >{\raggedright\arraybackslash}p{1.8cm}%
    >{\raggedright\arraybackslash\hyphenpenalty=10000}p{2.0cm}%
    >{\raggedright\arraybackslash}X%
    >{\raggedright\arraybackslash\hyphenpenalty=10000}p{2.8cm}@{}}
    % \caption*{Verzeichnis der verwendeten KI-Hilfsmittel}\\
    \toprule
    \textbf{Kennung} & \textbf{Modell} & \textbf{Datum} & \textbf{Prompt}
        & \textbf{Verwendung} \\
    \midrule
    \endfirsthead
    \toprule
    \textbf{Kennung} & \textbf{Modell} & \textbf{Datum} & \textbf{Prompt}
        & \textbf{Verwendung} \\
    \midrule
    \endhead
    \bottomrule
    \endlastfoot
    %
    Claude2026a & Anthropic Claude Opus 5.5 & 03.10.2026
      & Prüfe meine fertige Bachelorarbeit auf Fehler in Rechtschreibung, Zeichensetzung und Wortwahl und schlage für jede Stelle eine Korrektur vor. Inhalt, Aufbau und Argumentation bleiben unverändert. Ich entscheide bei jedem Vorschlag selbst, ob ich ihn übernehme.
      & gesamte Arbeit \\
    \addlinespace
    Claude2026b & Anthropic Claude Opus 5.5 & 03.10.2026
      & Schaue dir jeden Abschnitt an und gucke nach Schachtelsätzen, uneinheitlichen Fachbegriffen und schwerfälligen Formulierungen und schlage kürzere Alternativen vor. Halte dich dabei an meine beigefügten Stilregeln, etwa kurze Hauptsätze, ein Begriff je Sache und keine Verstärker.
      & gesamte Arbeit \\
    \addlinespace
    %
    Claude2026c & Anthropic Claude Opus 5 & 02.09.2026
      & Zu Kapitel 1 und 2 habe ich beim Gegenlesen eigene Anmerkungen notiert, sie hängen als Datei an. Ordne jede Anmerkung der betroffenen Textstelle zu und schlage je Stelle eine Änderung auf Satzebene vor.
      & Kap.~1, 2 \\
    \addlinespace
    Claude2026d & Anthropic Claude Opus 5 & 28.07.2026
      & Überarbeite den angehängten Absatz aus meinem Exposé für die Zielsetzung der Bachelorarbeit. Ziele und Gütemaße bleiben inhaltlich erhalten. Passe Begriffe und Formulierungen an meine angehängten Stichpunkte zum aktuellen Stand an.
      & Abschn.~1.2 \\
    \addlinespace
    Claude2026e & Anthropic Claude Opus 5 & 22.08.2026
      & In Shmuelis angehängtem Aufsatz geht es um den Unterschied zwischen erklärender und prognostischer Modellierung. Zieh mir die Definitionen beider Ansätze und die Stellen heraus, an denen der Aufsatz Folgen für die Modellwahl beschreibt, jeweils mit Seite.
      & Abschn.~2.2, 2.3, 2.7, 7.1 \\
    \addlinespace
    Claude2026f & Anthropic Claude Opus 5 & 28.08.2026
      & Suche mir frei zugängliche oder über die FOM erreichbare Standardwerke zu Grundlagen von Gütemaßen, Kreuzvalidierung und Hyperparametersuche. Nenne die Werke mit Kapitel oder Seite.
      & Abschn.~2.5 \\
    \addlinespace
    Claude2026g & Anthropic Claude Opus 5 & 02.09.2026
      & Lies im angehängten Text von Fahrmeir die Kapitel zu statistischen Tests und fasse zusammen, was dort zum Wilcoxon-Vorzeichen-Rang-Test, zur Gütefunktion und zum multiplen Testproblem steht. Ich brauche je Punkt die Seite.
      & Abschn.~2.5 \\
    \addlinespace
    Claude2026h & Anthropic Claude Opus 5 & 20.08.2026
      & Suche frei zugängliche Literatur zu Gütemaßen bei unausgewogenen Klassen, konkret zu Accuracy, Macro-F1 und mehrklassiger AUROC. Nenne je Werk die Seiten, auf denen das Thema behandelt wird, damit ich die Stellen nachlesen kann.
      & Abschn.~2.5, 6.2 \\
    \addlinespace
    Claude2026i & Anthropic Claude Opus 5 & 10.09.2026
      & Für jede der angehängten Studien brauche ich: Gegenstand, Datenbasis und zentralen Befund mit Seitenzahl. Formuliere den Befund in einem Satz.
      & Abschn.~2.7 \\
    \addlinespace
    Claude2026j & Anthropic Claude Opus 5 & 07.09.2026
      & Schlage mir den Rahmentext zur Quertabelle vor. Grundlage sind meine angehängten Stichpunkte zur Forschungslücke und zu den vier Punkten, in denen sich meine Arbeit von den bisherigen Studien unterscheidet. Der Text soll die Tabelle einordnen, nicht wiederholen.
      & Abschn.~2.7 \\
    \addlinespace
    Claude2026k & Anthropic Claude Opus 5 & 30.08.2026
      & Beschreibt das angehängte CRISP-DM-Handbuch von Chapman das Prozessmodell so, wie es der Aufsatz von Wirth und Hipp darstellt? Gleiche beide PDFs ab und nenne die Seiten, auf denen Chapman Phasen, Aufgaben und Abstraktionsebenen des Modells beschreibt.
      & Abschn.~3.1 \\
    \addlinespace
    Claude2026l & Anthropic Claude Opus 5 & 22.08.2026
      & Aus meinen angehängten Stichpunkten und der Datei der deskriptiven Auswertung soll ein Textvorschlag für die Datenbeschreibung und explorative Analyse entstehen. Übernimm nur Zahlen aus der Befunddatei, markiere jede Zahl, die du dort nicht findest.
      & Abschn.~4.2 \\
    \addlinespace
    Claude2026m & Anthropic Claude Opus 5 & 08.09.2026
      & Erkläre die beiden angehängten Codeausschnitte zur Verknüpfung der Zensusdaten mit den Stadtteilen und ordne ein, welche Rolle sie in der gesamten Datenaufbereitung spielen. Schlage mir dazu einen kurzen Text mit Bezug auf die Zeilennummern vor.
      & Abschn.~4.3 \\
    \addlinespace
    Claude2026n & Anthropic Claude Opus 5 & 09.09.2026
      & Fasse in Stichpunkten zusammen, was die angehängten Codeausschnitte zur Konstruktion der Merkmale tun und an welcher Stelle der Pipeline sie stehen. Nenne zu jedem Punkt die passenden Zeilennummern.
      & Abschn.~4.4 \\
    \addlinespace
    Claude2026o & Anthropic Claude Opus 5 & 09.09.2026
      & Welche Aufgabe haben die angehängten Codeausschnitte zu Fold-Zuteilung, Tests und Baselines im gesamten Ablauf der Auswertung? Mache mir einen kurzen Textvorschlag mit Bezug auf die Zeilennummern.
      & Abschn.~4.5 \\
    \addlinespace
    Claude2026p & Anthropic Claude Opus 5 & 25.08.2026
      & Im Anhang stehen die Verfahren und Tests, die für die Eignungsprüfung laufen: Ridge, Random Forest und XGBoost sowie RESET, Breusch-Pagan, Jarque-Bera und ein Überdispersionstest. Schreibe mir auf Basis der angehängten Arbeiten eine Beschreibung der Tests.
      & Abschn.~5.1 \\
    \addlinespace
    Claude2026q & Anthropic Claude Opus 5 & 07.09.2026
      & Schreibe mir zu jedem angehängten Codeausschnitt aus der Modellierung zwei bis drei Sätze: was der Code tut, wie er an die Datenaufbereitung anschließt und in welchen Zeilen das geschieht.
      & Abschn.~5.1, 5.2 \\
    \addlinespace
    Claude2026r & Anthropic Claude Opus 5 & 21.08.2026
      & Mache einen Vorschlag für eine Kürzung meines Ergebnisabschnitts um anderthalb Seiten, vor allem dort, wo der Text Tabellen nacherzählt oder Gesagtes wiederholt. Quellenverweise bleiben unangetastet, höchstens zwei Aussagen dürfen entfallen. Liste mir danach jede Streichung auf.
      & Abschn.~6.1, 6.2 \\
    \addlinespace
    Claude2026s & Anthropic Claude Opus 5 & 07.09.2026
      & Für den Abschnitt zum Aufwand fehlt mir die theoretische Laufzeitkomplexität der drei Verfahren. Schlage einen kurzen Textbaustein vor, der sich auf die angehängten Originalarbeiten stützt, und markiere jede Seitenangabe, die ich noch am Volltext prüfen muss.
      & Abschn.~6.4 \\
    \addlinespace
    Claude2026t & Anthropic Claude Opus 5 & 26.08.2026
      & Hier sind meine Stichpunkte zur Diskussion als TXT-Datei, dazu die Ergebnistabellen aus Kapitel~6. Schlage mir für Einordnung, Implikationen und Limitationen je einen Textvorschlag vor, der meine Gliederung übernimmt und nichts aus Kapitel~6 wiederholt.
      & Abschn.~7.1--7.3 \\
    \addlinespace
    Claude2026u & Anthropic Claude Opus 5 & 26.08.2026
      & Der angehängte Diskussionsentwurf unterschätzt aus meiner Sicht, was die Merkmale leisten. Prüfe Einordnung und Implikationen gegen meine Ergebnistabelle, ergänze einen Vergleich mit einer weiteren Referenz aus dieser Tabelle und schlage überarbeitete Absätze vor.
      & Abschn.~7.1, 7.2 \\
    \addlinespace
    Claude2026v & Anthropic Claude Opus 5 & 27.07.2026
      & Lege nach meinen angehängten Festlegungen eine zentrale Konfiguration für die Aufbereitung an: Analysezeitraum, Plausibilitätsfenster der Antwortzeit, ausgeschlossene Stadtteile und Schalter für die Downloads. Alle Skripte sollen diese Werte nur lesen und nie selbst setzen.
      & config.py \\
    \addlinespace
    Claude2026w & Anthropic Claude Opus 5 & 27.07.2026
      & Schreibe den Abruf der Rohdaten über die Schnittstellen von DataSF und dem Census Bureau. Welche Datensätze, Felder und Filter gebraucht werden, steht in meiner Quellenliste als Datei. Jede Quelle wird unverändert unter data/raw abgelegt.
      & s1\_\allowbreak{}daten.py \\
    \addlinespace
    Claude2026x & Anthropic Claude Opus 5 & 28.07.2026
      & Die Sozialdaten liegen je Census Tract vor. Ordne sie über die angehängte Zuordnungstabelle den Stadtteilen zu und gib je Jahrgang die Trefferquote aus.
      & s1\_\allowbreak{}daten.py, Code~3 \\
    \addlinespace
    Claude2026y & Anthropic Claude Opus 5 & 28.07.2026
      & Kaufman et~al. beschreiben Leakage auch durch Werte, die zum Vorhersagezeitpunkt noch gar nicht bekannt waren. Ordne deshalb jedem Monat nur den ACS-Jahrgang zu, der damals schon veröffentlicht war. Den Versatz je Jahrgang habe ich angehängt.
      & s1\_\allowbreak{}daten.py, Code~4 \\
    \addlinespace
    Claude2026z & Anthropic Claude Opus 5 & 28.07.2026
      & Bilde die Anteilswerte aus meiner angehängten Merkmalsliste. Ist ein Nenner null, soll ein fehlender Wert entstehen.
      & s1\_\allowbreak{}daten.py, Code~6 \\
    \addlinespace
    Claude2026aa & Anthropic Claude Opus 5 & 29.07.2026
      & Für die Kriminalität habe ich mich nach Brantingham und Brantingham für einen Standortquotienten entschieden. Implementiere ihn als relativen Index je Stadtteil und Monat, bezogen auf die gesamte Stadt im selben Monat.
      & s1\_\allowbreak{}daten.py, Code~7 \\
    \addlinespace
    Claude2026ab & Anthropic Claude Opus 5 & 29.07.2026
      & Aggregiere die Einsätze auf Stadtteil und Monat, einmal als Anzahl und einmal als überwiegende Einsatzart. Jede Kombination soll als Zeile vorhanden sein, jährliche Werte werden nach meiner angehängten Regel vorwärts gefüllt.
      & s2\_\allowbreak{}datensaetze.py, Code~5 \\
    \addlinespace
    Claude2026ac & Anthropic Claude Opus 5 & 29.07.2026
      & Kodiere den Kalendermonat als Sinus und Kosinus und ergänze die verzögerten Werte aus meiner Merkmalsliste, nur aus abgeschlossenen Vormonaten.
      & s2\_\allowbreak{}datensaetze.py, Code~8 \\
    \addlinespace
    Claude2026ad & Anthropic Claude Opus 5 & 30.07.2026
      & Roberts et~al. empfehlen bei räumlich abhängigen Daten eine gruppierte Kreuzvalidierung. Teile deshalb ganze Stadtteile auf fünf Folds und ein Hold-out auf, geschichtet nach brand-dominierten Monaten und Bevölkerung. Die Zuteilung steht als Spalte im Datensatz.
      & s2\_\allowbreak{}datensaetze.py, Code~10 \\
    \addlinespace
    Claude2026ae & Anthropic Claude Opus 5 & 30.07.2026
      & Baue eine Prüfung ein, die nach dem Bau des Panels abbricht, wenn Stadtteile, Zeitraum oder Wohnbevölkerung nicht zu meinen angehängten Sollwerten passen. Die Fehlermeldung soll den betroffenen Stadtteil und den gemessenen Wert nennen.
      & s2\_\allowbreak{}datensaetze.py, Code~1 \\
    \addlinespace
    Claude2026af & Anthropic Claude Opus 5 & 31.07.2026
      & Wie mit meinem Betreuer abgestimmt, messen sich die Verfahren an zwei Baselines. Implementiere eine Referenz ohne Merkmale und ein Poisson-GLM mit der Bevölkerung als Offset nach Winkelmann, für die Einsatzart ein multinomiales Logit.
      & v1\_\allowbreak{}baselines.py, Code~12, 13 \\
    \addlinespace
    Claude2026ag & Anthropic Claude Opus 5 & 03.08.2026
      & Setze die angehängten Eignungsprüfungen, die ich vor der Modellierung festgelegt habe, als ein Skript mit einem Bericht am Ende um. Dazu gehören die Überdispersion als Hilfsregression, die Linearität grafisch vor dem Einsatz von Ridge, der RESET nach Ramsey und je Verfahren eine Tabelle mit Anforderung, Teststatistik und p-Wert.
      & v2\_\allowbreak{}eignung.py, Code~14 \\
    \addlinespace
    Claude2026ah & Anthropic Claude Opus 5 & 03.08.2026
      & Fasse die Aufbereitung in einem einzigen Befehl zusammen, der beide Schritte nacheinander ausführt, bei einem Fehler abbricht und am Ende Zeilen und Spalten beider Dateien ausgibt. Dasselbe brauche ich für die Vorprüfung: erst die Baselines, danach die Eignungsprüfung, die deren Werte nur liest.
      & build.py, run.py \\
    \addlinespace
    Claude2026ai & Anthropic Claude Opus 5 & 04.08.2026
      & Lege das Skript für den Mengenvergleich in Phasen an, die ich nacheinander fülle: Tuning, Bewertung, Aggregation und Vergleich. Ridge rechnet nach meiner Methodik auf $\log(1+y)$, Random Forest und XGBoost mit zählgerechter Verlustfunktion. Ein einzelner Lauf passt auf der Rate an, rechnet auf Einsätze zurück und misst die Zeit.
      & m02\_\allowbreak{}menge.py, Code~15, 16 \\
    \addlinespace
    Claude2026aj & Anthropic Claude Opus 5 & 05.08.2026
      & Übertrage den Aufbau des Mengenvergleichs auf die Einsatzart mit Random Forest und XGBoost. Die vier Klassen sollen global einheitlich kodiert und ausgeglichen gewichtet werden, Hauptmaß ist nach Opitz und Burst der Macro-F1.
      & m03\_\allowbreak{}struktur.py, Code~17 \\
    \addlinespace
    Claude2026ak & Anthropic Claude Opus 5 & 06.08.2026
      & Die zehn Wiederholungen verschieben bisher nur die Nummern der Folds, die Zusammensetzung bleibt gleich, und das Hold-out wandert mit. Korrigiere das so, dass innerhalb der Rangblöcke meiner Schichtung neu gemischt wird und das Hold-out fest bleibt.
      & v0\_\allowbreak{}aufteilung.py \\
    \addlinespace
    Claude2026al & Anthropic Claude Opus 5 & 06.08.2026
      & Nach Bergstra und Bengio nehme ich eine Zufallssuche. Setze sie als innere, nach Stadtteilen gruppierte Kreuzvalidierung mit meinem angehängten Suchraum um.
      & m02\_\allowbreak{}menge.py, Code~18 \\
    \addlinespace
    Claude2026am & Anthropic Claude Opus 5 & 07.08.2026
      & Hurlbert beschreibt Pseudoreplikation, wenn abhängige Messungen als unabhängig gezählt werden. Mittle die Gütemaße deshalb zweistufig, erst über die Folds je Wiederholung und dann über die zehn Wiederholungen, und gib beide Streuungen aus.
      & m02\_\allowbreak{}menge.py, Code~11, 19 \\
    \addlinespace
    Claude2026an & Anthropic Claude Opus 5 & 07.08.2026
      & Die Eignungsprüfung verwirft die lineare Form nur innerhalb der Stichprobe. Prüfe in einem eigenen Skript, ob Quadrat- und Interaktionsterme im Poisson-GLM auch auf unbekannten Stadtteilen besser vorhersagen, mit denselben Folds wie im Vergleich.
      & v3\_\allowbreak{}spezifikation.py \\
    \addlinespace
    Claude2026ao & Anthropic Claude Opus 5 & 08.08.2026
      & Vergleiche jedes Verfahren gepaart mit dem Wilcoxon-Test gegen die Stufe-2-Baseline und korrigiere nach Holm. Die Schlussbewertung auf dem Hold-out läuft nur mit ausdrücklichem Argument, einmalig und mit den Parametern aus dem Tuning.
      & m02\_\allowbreak{}menge.py, m03\_\allowbreak{}struktur.py \\
    \addlinespace
    Claude2026ap & Anthropic Claude Opus 5 & 08.08.2026
      & Wiederholung null dient zugleich dem Tuning und der Bewertung. Beziffere, wie viel Vorsprung die Verfahren dadurch bekommen.
      & m02\_\allowbreak{}menge.py, m03\_\allowbreak{}struktur.py \\
    \addlinespace
    Claude2026aq & Anthropic Claude Opus 5 & 11.08.2026
      & Ich möchte die fertigen Dateien prüfen und nicht den Code. Schreibe dazu eine eigene Testdatei nach meiner angehängten Liste, gegliedert nach Panel, Aufteilung, Merkmalen und Struktur. Darunter fallen keine Ergebnisvariable als Merkmal im Sinne von Kaufman et~al. und kein Stadtteil zugleich in Training und Test. Sie soll auch ohne pytest laufen.
      & test\_\allowbreak{}aufbereitung.py, Code~2, 9 \\
    \addlinespace
    Claude2026ar & Anthropic Claude Opus 5 & 13.08.2026
      & SHAP-Werte nach Lundberg und Lee zeigen nur, wie ein Modell seine Aufmerksamkeit verteilt. Ergänze deshalb eine Ablation je Faktorgruppe aus meiner Gruppierung, bei der das Modell ohne die Gruppe neu angepasst und bewertet wird.
      & m04\_\allowbreak{}shap.py \\
    \addlinespace
    Claude2026as & Anthropic Claude Opus 5 & 14.08.2026
      & Sorge dafür, dass ein zweiter Lauf dieselben Zahlen liefert: ein fester Startwert für alle Zufallsschritte, Bewertung auf einem Kern und eine Datei mit den genauen Paketversionen. Die Konstanten der Modellierung kommen in eine eigene Konfiguration.
      & config\_\allowbreak{}modelle.py, requirements\_\allowbreak{}lauf.txt \\
    \addlinespace
    Claude2026at & Anthropic Claude Opus 5 & 15.08.2026
      & Ob das Budget der Hyperparametersuche reicht, will ich wie von Bischl et~al. empfohlen prüfen. Wiederhole die Suche mit steigender Zahl an Versuchen, zeige den Verlauf als Kurve und ergänze einen kurzen Probelauf.
      & suchdiagnose.py \\
    \addlinespace
    Claude2026au & Anthropic Claude Opus 5 & 17.08.2026
      & Ein Macro-F1 lässt sich ohne Obergrenze schlecht einordnen. Schätze mit meinem angehängten Ansatz, wie gut die Einsatzart mit diesen Merkmalen höchstens vorhersagbar ist, und setze die Werte aus dem Vergleich ins Verhältnis dazu.
      & v4\_\allowbreak{}decke.py \\
    \addlinespace
    Claude2026av & Anthropic Claude Opus 5 & 20.08.2026
      & Zum Data Understanding nach Chapman et~al. brauche ich zwei getrennte Skripte: eine Merkmalstabelle mit Skalenniveau, Einheit und Quelle aus meiner Vorlage und eine deskriptive Auswertung mit Verteilung, Varianzzerlegung und Korrelationen.
      & codebook.py, deskriptiv.py \\
    \addlinespace
    Claude2026aw & Anthropic Claude Opus 5 & 23.08.2026
      & Der Census-Schlüssel steht im Klartext in der Konfiguration. Lies ihn stattdessen aus einer Umgebungsvariable.
      & config.py \\
    \addlinespace
    Claude2026ax & Anthropic Claude Opus 5 & 25.08.2026
      & Stelle aus den Rohquellen die Befunde zusammen, aus denen ein Eingriff in der Aufbereitung folgt, etwa Dubletten, fehlende Baujahre oder Lücken in der Abdeckung. Meine Liste der Eingriffe liegt bei, alles andere bleibt draußen.
      & rohbefunde.py \\
    \addlinespace
    Claude2026ay & Anthropic Claude Opus 5 & 29.08.2026
      & In den ACS-Jahrgängen mit älteren Zensusgrenzen fehlt ein Teil der Wohnbevölkerung. Suche die Ursache in der Zuordnung der Tracts, schlage eine Korrektur vor und lass die Aufbereitung abbrechen, wenn zu wenig Bevölkerung zugeordnet ist.
      & s1\_\allowbreak{}daten.py, Code~3 \\
    \addlinespace
    Claude2026az & Anthropic Claude Opus 5 & 30.08.2026
      & Bevor ich Ergebnisse ansehe, will ich wissen, welche Stadtteile im Hold-out und welche in der Entwicklung liegen. Erstelle ein Profil beider Hälften mit Klassenverteilung und Zielgrößen, unabhängig von jedem Modellergebnis.
      & panelprofil.py \\
    \addlinespace
    Claude2026ba & Anthropic Claude Opus 5 & 02.09.2026
      & Probst et~al. zeigen, dass Verfahren unterschiedlich stark auf ihre Hyperparameter reagieren. Setze jeden getunten Parametersatz auf allen anderen Folds ein und fasse zusammen, wie weit die Güte dabei schwankt.
      & parameter\allowbreak{}sensitivitaet.py \\
    \addlinespace
    Claude2026bb & Anthropic Claude Opus 5 & 11.09.2026
      & Nach Altman und Bland ist ein nicht signifikanter Test kein Beleg für Gleichheit. Berechne für die gepaarten Vergleiche Effektstärke und Trennschärfe bei zehn Wiederholungsmitteln, mit derselben Formel für die Streuung wie in den Vergleichsskripten.
      & trennschaerfe.py \\
    \addlinespace
    Claude2026bc & Anthropic Claude Opus 5 & 14.09.2026
      & Setze die README für die Abgabe auf, gegliedert nach Umgebung, Daten, Ablauf, Reproduzierbarkeit und Aufbau. Den Ablauf gebe ich als Befehlsfolge in einer Datei vor. Alle weiteren Inhalte stammen aus Code und Konfiguration, dazu ein Hinweis, dass ein neuer Abruf die Rohdaten überschreibt und der Census-Schlüssel nicht beiliegt.
      & README.md \\
    \addlinespace
    Claude2026bd & Anthropic Claude Opus 5.5 & 24.09.2026
      & Gleiche die Kommentare aller Skripte mit dem Stand der Arbeit ab. Meine angehängte Liste nennt veraltete Begriffe, Zahlen und Kapitelverweise. Schlage je Datei die Änderungen vor, geändert werden nur Kommentare und Docstrings.
      & alle Skripte \\
    \addlinespace
    Claude2026be & Anthropic Claude Opus 5.5 & 24.09.2026
      & Bringe die Kommentare in allen Dateien in eine einheitliche Form nach meiner Vorlage: ein Kopf mit Zweck, Aufruf, Eingang und Ausgang, darunter kurze Stichpunkte zu den Entscheidungen. Interne Kürzel und Verweise auf Arbeitsnotizen entfallen.
      & alle Skripte \\
    \addlinespace
\end{xltabular}
% \tabquelle
\endgroup

% Benoetigt in der Praeambel: \usepackage{xltabular}
% Regeln: Projekt-Notiz KI-Verzeichnis_Regeln_2026-09-30
%   Spalten: Kennung, Modell, Datum, Prompt, Verwendung
%   Kennung: Claude2026a ... Claude2026z, Claude2026aa, Claude2026ab ...
%   Sortierung: nach Vorkommen, erst Gesamttext, dann Abschnitte
%               und Literatur, dann Code in der Reihenfolge der Entstehung
%               (Pipeline: Aufbereitung, Vorpruefung, Modelle, README,
%               Kommentare)
%   Prompts ohne Anfuehrungszeichen, ohne Kursivschrift
% ---------------------------------------------------
\newpage
\phantomsection
\section*{KI-Verzeichnis}
\addcontentsline{toc}{section}{KI-Verzeichnis}

\begingroup
\tabformat
\setlength{\tabcolsep}{4pt}
\begin{xltabular}{\textwidth}{@{}%
    >{\raggedright\arraybackslash\hyphenpenalty=10000}p{2.6cm}%
    >{\raggedright\arraybackslash}p{1.8cm}%
    >{\raggedright\arraybackslash\hyphenpenalty=10000}p{2.0cm}%
    >{\raggedright\arraybackslash}X%
    >{\raggedright\arraybackslash\hyphenpenalty=10000}p{2.8cm}@{}}
    % \caption*{Verzeichnis der verwendeten KI-Hilfsmittel}\\
    \toprule
    \textbf{Kennung} & \textbf{Modell} & \textbf{Datum} & \textbf{Prompt}
        & \textbf{Verwendung} \\
    \midrule
    \endfirsthead
    \toprule
    \textbf{Kennung} & \textbf{Modell} & \textbf{Datum} & \textbf{Prompt}
        & \textbf{Verwendung} \\
    \midrule
    \endhead
    \bottomrule
    \endlastfoot
    %
    Claude2026a & Anthropic Claude Opus 5.5 & 03.10.2026
      & Prüfe meine fertige Bachelorarbeit auf Fehler in Rechtschreibung, Zeichensetzung und Wortwahl und schlage für jede Stelle eine Korrektur vor. Inhalt, Aufbau und Argumentation bleiben unverändert. Ich entscheide bei jedem Vorschlag selbst, ob ich ihn übernehme.
      & gesamte Arbeit \\
    \addlinespace
    Claude2026b & Anthropic Claude Opus 5.5 & 03.10.2026
      & Schaue dir jeden Abschnitt an und gucke nach Schachtelsätzen, uneinheitlichen Fachbegriffen und schwerfälligen Formulierungen und schlage kürzere Alternativen vor. Halte dich dabei an meine beigefügten Stilregeln, etwa kurze Hauptsätze, ein Begriff je Sache und keine Verstärker.
      & gesamte Arbeit \\
    \addlinespace
    %
    Claude2026c & Anthropic Claude Opus 5 & 02.09.2026
      & Zu Kapitel 1 und 2 habe ich beim Gegenlesen eigene Anmerkungen notiert, sie hängen als Datei an. Ordne jede Anmerkung der betroffenen Textstelle zu und schlage je Stelle eine Änderung auf Satzebene vor.
      & Kap.~1, 2 \\
    \addlinespace
    Claude2026d & Anthropic Claude Opus 5 & 28.07.2026
      & Überarbeite den angehängten Absatz aus meinem Exposé für die Zielsetzung der Bachelorarbeit. Ziele und Gütemaße bleiben inhaltlich erhalten. Passe Begriffe und Formulierungen an meine angehängten Stichpunkte zum aktuellen Stand an.
      & Abschn.~1.2 \\
    \addlinespace
    Claude2026e & Anthropic Claude Opus 5 & 22.08.2026
      & In Shmuelis angehängtem Aufsatz geht es um den Unterschied zwischen erklärender und prognostischer Modellierung. Zieh mir die Definitionen beider Ansätze und die Stellen heraus, an denen der Aufsatz Folgen für die Modellwahl beschreibt, jeweils mit Seite.
      & Abschn.~2.2, 2.3, 2.7, 7.1 \\
    \addlinespace
    Claude2026f & Anthropic Claude Opus 5 & 28.08.2026
      & Suche mir frei zugängliche oder über die FOM erreichbare Standardwerke zu Grundlagen von Gütemaßen, Kreuzvalidierung und Hyperparametersuche. Nenne die Werke mit Kapitel oder Seite.
      & Abschn.~2.5 \\
    \addlinespace
    Claude2026g & Anthropic Claude Opus 5 & 02.09.2026
      & Lies im angehängten Text von Fahrmeir die Kapitel zu statistischen Tests und fasse zusammen, was dort zum Wilcoxon-Vorzeichen-Rang-Test, zur Gütefunktion und zum multiplen Testproblem steht. Ich brauche je Punkt die Seite.
      & Abschn.~2.5 \\
    \addlinespace
    Claude2026h & Anthropic Claude Opus 5 & 20.08.2026
      & Suche frei zugängliche Literatur zu Gütemaßen bei unausgewogenen Klassen, konkret zu Accuracy, Macro-F1 und mehrklassiger AUROC. Nenne je Werk die Seiten, auf denen das Thema behandelt wird, damit ich die Stellen nachlesen kann.
      & Abschn.~2.5, 6.2 \\
    \addlinespace
    Claude2026i & Anthropic Claude Opus 5 & 10.09.2026
      & Für jede der angehängten Studien brauche ich: Gegenstand, Datenbasis und zentralen Befund mit Seitenzahl. Formuliere den Befund in einem Satz.
      & Abschn.~2.7 \\
    \addlinespace
    Claude2026j & Anthropic Claude Opus 5 & 07.09.2026
      & Schlage mir den Rahmentext zur Quertabelle vor. Grundlage sind meine angehängten Stichpunkte zur Forschungslücke und zu den vier Punkten, in denen sich meine Arbeit von den bisherigen Studien unterscheidet. Der Text soll die Tabelle einordnen, nicht wiederholen.
      & Abschn.~2.7 \\
    \addlinespace
    Claude2026k & Anthropic Claude Opus 5 & 30.08.2026
      & Beschreibt das angehängte CRISP-DM-Handbuch von Chapman das Prozessmodell so, wie es der Aufsatz von Wirth und Hipp darstellt? Gleiche beide PDFs ab und nenne die Seiten, auf denen Chapman Phasen, Aufgaben und Abstraktionsebenen des Modells beschreibt.
      & Abschn.~3.1 \\
    \addlinespace
    Claude2026l & Anthropic Claude Opus 5 & 22.08.2026
      & Aus meinen angehängten Stichpunkten und der Datei der deskriptiven Auswertung soll ein Textvorschlag für die Datenbeschreibung und explorative Analyse entstehen. Übernimm nur Zahlen aus der Befunddatei, markiere jede Zahl, die du dort nicht findest.
      & Abschn.~4.2 \\
    \addlinespace
    Claude2026m & Anthropic Claude Opus 5 & 08.09.2026
      & Erkläre die beiden angehängten Codeausschnitte zur Verknüpfung der Zensusdaten mit den Stadtteilen und ordne ein, welche Rolle sie in der gesamten Datenaufbereitung spielen. Schlage mir dazu einen kurzen Text mit Bezug auf die Zeilennummern vor.
      & Abschn.~4.3 \\
    \addlinespace
    Claude2026n & Anthropic Claude Opus 5 & 09.09.2026
      & Fasse in Stichpunkten zusammen, was die angehängten Codeausschnitte zur Konstruktion der Merkmale tun und an welcher Stelle der Pipeline sie stehen. Nenne zu jedem Punkt die passenden Zeilennummern.
      & Abschn.~4.4 \\
    \addlinespace
    Claude2026o & Anthropic Claude Opus 5 & 09.09.2026
      & Welche Aufgabe haben die angehängten Codeausschnitte zu Fold-Zuteilung, Tests und Baselines im gesamten Ablauf der Auswertung? Mache mir einen kurzen Textvorschlag mit Bezug auf die Zeilennummern.
      & Abschn.~4.5 \\
    \addlinespace
    Claude2026p & Anthropic Claude Opus 5 & 25.08.2026
      & Im Anhang stehen die Verfahren und Tests, die für die Eignungsprüfung laufen: Ridge, Random Forest und XGBoost sowie RESET, Breusch-Pagan, Jarque-Bera und ein Überdispersionstest. Schreibe mir auf Basis der angehängten Arbeiten eine Beschreibung der Tests.
      & Abschn.~5.1 \\
    \addlinespace
    Claude2026q & Anthropic Claude Opus 5 & 07.09.2026
      & Schreibe mir zu jedem angehängten Codeausschnitt aus der Modellierung zwei bis drei Sätze: was der Code tut, wie er an die Datenaufbereitung anschließt und in welchen Zeilen das geschieht.
      & Abschn.~5.1, 5.2 \\
    \addlinespace
    Claude2026r & Anthropic Claude Opus 5 & 21.08.2026
      & Mache einen Vorschlag für eine Kürzung meines Ergebnisabschnitts um anderthalb Seiten, vor allem dort, wo der Text Tabellen nacherzählt oder Gesagtes wiederholt. Quellenverweise bleiben unangetastet, höchstens zwei Aussagen dürfen entfallen. Liste mir danach jede Streichung auf.
      & Abschn.~6.1, 6.2 \\
    \addlinespace
    Claude2026s & Anthropic Claude Opus 5 & 07.09.2026
      & Für den Abschnitt zum Aufwand fehlt mir die theoretische Laufzeitkomplexität der drei Verfahren. Schlage einen kurzen Textbaustein vor, der sich auf die angehängten Originalarbeiten stützt, und markiere jede Seitenangabe, die ich noch am Volltext prüfen muss.
      & Abschn.~6.4 \\
    \addlinespace
    Claude2026t & Anthropic Claude Opus 5 & 26.08.2026
      & Hier sind meine Stichpunkte zur Diskussion als TXT-Datei, dazu die Ergebnistabellen aus Kapitel~6. Schlage mir für Einordnung, Implikationen und Limitationen je einen Textvorschlag vor, der meine Gliederung übernimmt und nichts aus Kapitel~6 wiederholt.
      & Abschn.~7.1--7.3 \\
    \addlinespace
    Claude2026u & Anthropic Claude Opus 5 & 26.08.2026
      & Der angehängte Diskussionsentwurf unterschätzt aus meiner Sicht, was die Merkmale leisten. Prüfe Einordnung und Implikationen gegen meine Ergebnistabelle, ergänze einen Vergleich mit einer weiteren Referenz aus dieser Tabelle und schlage überarbeitete Absätze vor.
      & Abschn.~7.1, 7.2 \\
    \addlinespace
    Claude2026v & Anthropic Claude Opus 5 & 27.07.2026
      & Lege nach meinen angehängten Festlegungen eine zentrale Konfiguration für die Aufbereitung an: Analysezeitraum, Plausibilitätsfenster der Antwortzeit, ausgeschlossene Stadtteile und Schalter für die Downloads. Alle Skripte sollen diese Werte nur lesen und nie selbst setzen.
      & config.py \\
    \addlinespace
    Claude2026w & Anthropic Claude Opus 5 & 27.07.2026
      & Schreibe den Abruf der Rohdaten über die Schnittstellen von DataSF und dem Census Bureau. Welche Datensätze, Felder und Filter gebraucht werden, steht in meiner Quellenliste als Datei. Jede Quelle wird unverändert unter data/raw abgelegt.
      & s1\_\allowbreak{}daten.py \\
    \addlinespace
    Claude2026x & Anthropic Claude Opus 5 & 28.07.2026
      & Die Sozialdaten liegen je Census Tract vor. Ordne sie über die angehängte Zuordnungstabelle den Stadtteilen zu und gib je Jahrgang die Trefferquote aus.
      & s1\_\allowbreak{}daten.py, Code~3 \\
    \addlinespace
    Claude2026y & Anthropic Claude Opus 5 & 28.07.2026
      & Kaufman et~al. beschreiben Leakage auch durch Werte, die zum Vorhersagezeitpunkt noch gar nicht bekannt waren. Ordne deshalb jedem Monat nur den ACS-Jahrgang zu, der damals schon veröffentlicht war. Den Versatz je Jahrgang habe ich angehängt.
      & s1\_\allowbreak{}daten.py, Code~4 \\
    \addlinespace
    Claude2026z & Anthropic Claude Opus 5 & 28.07.2026
      & Bilde die Anteilswerte aus meiner angehängten Merkmalsliste. Ist ein Nenner null, soll ein fehlender Wert entstehen.
      & s1\_\allowbreak{}daten.py, Code~6 \\
    \addlinespace
    Claude2026aa & Anthropic Claude Opus 5 & 29.07.2026
      & Für die Kriminalität habe ich mich nach Brantingham und Brantingham für einen Standortquotienten entschieden. Implementiere ihn als relativen Index je Stadtteil und Monat, bezogen auf die gesamte Stadt im selben Monat.
      & s1\_\allowbreak{}daten.py, Code~7 \\
    \addlinespace
    Claude2026ab & Anthropic Claude Opus 5 & 29.07.2026
      & Aggregiere die Einsätze auf Stadtteil und Monat, einmal als Anzahl und einmal als überwiegende Einsatzart. Jede Kombination soll als Zeile vorhanden sein, jährliche Werte werden nach meiner angehängten Regel vorwärts gefüllt.
      & s2\_\allowbreak{}datensaetze.py, Code~5 \\
    \addlinespace
    Claude2026ac & Anthropic Claude Opus 5 & 29.07.2026
      & Kodiere den Kalendermonat als Sinus und Kosinus und ergänze die verzögerten Werte aus meiner Merkmalsliste, nur aus abgeschlossenen Vormonaten.
      & s2\_\allowbreak{}datensaetze.py, Code~8 \\
    \addlinespace
    Claude2026ad & Anthropic Claude Opus 5 & 30.07.2026
      & Roberts et~al. empfehlen bei räumlich abhängigen Daten eine gruppierte Kreuzvalidierung. Teile deshalb ganze Stadtteile auf fünf Folds und ein Hold-out auf, geschichtet nach brand-dominierten Monaten und Bevölkerung. Die Zuteilung steht als Spalte im Datensatz.
      & s2\_\allowbreak{}datensaetze.py, Code~10 \\
    \addlinespace
    Claude2026ae & Anthropic Claude Opus 5 & 30.07.2026
      & Baue eine Prüfung ein, die nach dem Bau des Panels abbricht, wenn Stadtteile, Zeitraum oder Wohnbevölkerung nicht zu meinen angehängten Sollwerten passen. Die Fehlermeldung soll den betroffenen Stadtteil und den gemessenen Wert nennen.
      & s2\_\allowbreak{}datensaetze.py, Code~1 \\
    \addlinespace
    Claude2026af & Anthropic Claude Opus 5 & 31.07.2026
      & Wie mit meinem Betreuer abgestimmt, messen sich die Verfahren an zwei Baselines. Implementiere eine Referenz ohne Merkmale und ein Poisson-GLM mit der Bevölkerung als Offset nach Winkelmann, für die Einsatzart ein multinomiales Logit.
      & v1\_\allowbreak{}baselines.py, Code~12, 13 \\
    \addlinespace
    Claude2026ag & Anthropic Claude Opus 5 & 03.08.2026
      & Setze die angehängten Eignungsprüfungen, die ich vor der Modellierung festgelegt habe, als ein Skript mit einem Bericht am Ende um. Dazu gehören die Überdispersion als Hilfsregression, die Linearität grafisch vor dem Einsatz von Ridge, der RESET nach Ramsey und je Verfahren eine Tabelle mit Anforderung, Teststatistik und p-Wert.
      & v2\_\allowbreak{}eignung.py, Code~14 \\
    \addlinespace
    Claude2026ah & Anthropic Claude Opus 5 & 03.08.2026
      & Fasse die Aufbereitung in einem einzigen Befehl zusammen, der beide Schritte nacheinander ausführt, bei einem Fehler abbricht und am Ende Zeilen und Spalten beider Dateien ausgibt. Dasselbe brauche ich für die Vorprüfung: erst die Baselines, danach die Eignungsprüfung, die deren Werte nur liest.
      & build.py, run.py \\
    \addlinespace
    Claude2026ai & Anthropic Claude Opus 5 & 04.08.2026
      & Lege das Skript für den Mengenvergleich in Phasen an, die ich nacheinander fülle: Tuning, Bewertung, Aggregation und Vergleich. Ridge rechnet nach meiner Methodik auf $\log(1+y)$, Random Forest und XGBoost mit zählgerechter Verlustfunktion. Ein einzelner Lauf passt auf der Rate an, rechnet auf Einsätze zurück und misst die Zeit.
      & m02\_\allowbreak{}menge.py, Code~15, 16 \\
    \addlinespace
    Claude2026aj & Anthropic Claude Opus 5 & 05.08.2026
      & Übertrage den Aufbau des Mengenvergleichs auf die Einsatzart mit Random Forest und XGBoost. Die vier Klassen sollen global einheitlich kodiert und ausgeglichen gewichtet werden, Hauptmaß ist nach Opitz und Burst der Macro-F1.
      & m03\_\allowbreak{}struktur.py, Code~17 \\
    \addlinespace
    Claude2026ak & Anthropic Claude Opus 5 & 06.08.2026
      & Die zehn Wiederholungen verschieben bisher nur die Nummern der Folds, die Zusammensetzung bleibt gleich, und das Hold-out wandert mit. Korrigiere das so, dass innerhalb der Rangblöcke meiner Schichtung neu gemischt wird und das Hold-out fest bleibt.
      & v0\_\allowbreak{}aufteilung.py \\
    \addlinespace
    Claude2026al & Anthropic Claude Opus 5 & 06.08.2026
      & Nach Bergstra und Bengio nehme ich eine Zufallssuche. Setze sie als innere, nach Stadtteilen gruppierte Kreuzvalidierung mit meinem angehängten Suchraum um.
      & m02\_\allowbreak{}menge.py, Code~18 \\
    \addlinespace
    Claude2026am & Anthropic Claude Opus 5 & 07.08.2026
      & Hurlbert beschreibt Pseudoreplikation, wenn abhängige Messungen als unabhängig gezählt werden. Mittle die Gütemaße deshalb zweistufig, erst über die Folds je Wiederholung und dann über die zehn Wiederholungen, und gib beide Streuungen aus.
      & m02\_\allowbreak{}menge.py, Code~11, 19 \\
    \addlinespace
    Claude2026an & Anthropic Claude Opus 5 & 07.08.2026
      & Die Eignungsprüfung verwirft die lineare Form nur innerhalb der Stichprobe. Prüfe in einem eigenen Skript, ob Quadrat- und Interaktionsterme im Poisson-GLM auch auf unbekannten Stadtteilen besser vorhersagen, mit denselben Folds wie im Vergleich.
      & v3\_\allowbreak{}spezifikation.py \\
    \addlinespace
    Claude2026ao & Anthropic Claude Opus 5 & 08.08.2026
      & Vergleiche jedes Verfahren gepaart mit dem Wilcoxon-Test gegen die Stufe-2-Baseline und korrigiere nach Holm. Die Schlussbewertung auf dem Hold-out läuft nur mit ausdrücklichem Argument, einmalig und mit den Parametern aus dem Tuning.
      & m02\_\allowbreak{}menge.py, m03\_\allowbreak{}struktur.py \\
    \addlinespace
    Claude2026ap & Anthropic Claude Opus 5 & 08.08.2026
      & Wiederholung null dient zugleich dem Tuning und der Bewertung. Beziffere, wie viel Vorsprung die Verfahren dadurch bekommen.
      & m02\_\allowbreak{}menge.py, m03\_\allowbreak{}struktur.py \\
    \addlinespace
    Claude2026aq & Anthropic Claude Opus 5 & 11.08.2026
      & Ich möchte die fertigen Dateien prüfen und nicht den Code. Schreibe dazu eine eigene Testdatei nach meiner angehängten Liste, gegliedert nach Panel, Aufteilung, Merkmalen und Struktur. Darunter fallen keine Ergebnisvariable als Merkmal im Sinne von Kaufman et~al. und kein Stadtteil zugleich in Training und Test. Sie soll auch ohne pytest laufen.
      & test\_\allowbreak{}aufbereitung.py, Code~2, 9 \\
    \addlinespace
    Claude2026ar & Anthropic Claude Opus 5 & 13.08.2026
      & SHAP-Werte nach Lundberg und Lee zeigen nur, wie ein Modell seine Aufmerksamkeit verteilt. Ergänze deshalb eine Ablation je Faktorgruppe aus meiner Gruppierung, bei der das Modell ohne die Gruppe neu angepasst und bewertet wird.
      & m04\_\allowbreak{}shap.py \\
    \addlinespace
    Claude2026as & Anthropic Claude Opus 5 & 14.08.2026
      & Sorge dafür, dass ein zweiter Lauf dieselben Zahlen liefert: ein fester Startwert für alle Zufallsschritte, Bewertung auf einem Kern und eine Datei mit den genauen Paketversionen. Die Konstanten der Modellierung kommen in eine eigene Konfiguration.
      & config\_\allowbreak{}modelle.py, requirements\_\allowbreak{}lauf.txt \\
    \addlinespace
    Claude2026at & Anthropic Claude Opus 5 & 15.08.2026
      & Ob das Budget der Hyperparametersuche reicht, will ich wie von Bischl et~al. empfohlen prüfen. Wiederhole die Suche mit steigender Zahl an Versuchen, zeige den Verlauf als Kurve und ergänze einen kurzen Probelauf.
      & suchdiagnose.py \\
    \addlinespace
    Claude2026au & Anthropic Claude Opus 5 & 17.08.2026
      & Ein Macro-F1 lässt sich ohne Obergrenze schlecht einordnen. Schätze mit meinem angehängten Ansatz, wie gut die Einsatzart mit diesen Merkmalen höchstens vorhersagbar ist, und setze die Werte aus dem Vergleich ins Verhältnis dazu.
      & v4\_\allowbreak{}decke.py \\
    \addlinespace
    Claude2026av & Anthropic Claude Opus 5 & 20.08.2026
      & Zum Data Understanding nach Chapman et~al. brauche ich zwei getrennte Skripte: eine Merkmalstabelle mit Skalenniveau, Einheit und Quelle aus meiner Vorlage und eine deskriptive Auswertung mit Verteilung, Varianzzerlegung und Korrelationen.
      & codebook.py, deskriptiv.py \\
    \addlinespace
    Claude2026aw & Anthropic Claude Opus 5 & 23.08.2026
      & Der Census-Schlüssel steht im Klartext in der Konfiguration. Lies ihn stattdessen aus einer Umgebungsvariable.
      & config.py \\
    \addlinespace
    Claude2026ax & Anthropic Claude Opus 5 & 25.08.2026
      & Stelle aus den Rohquellen die Befunde zusammen, aus denen ein Eingriff in der Aufbereitung folgt, etwa Dubletten, fehlende Baujahre oder Lücken in der Abdeckung. Meine Liste der Eingriffe liegt bei, alles andere bleibt draußen.
      & rohbefunde.py \\
    \addlinespace
    Claude2026ay & Anthropic Claude Opus 5 & 29.08.2026
      & In den ACS-Jahrgängen mit älteren Zensusgrenzen fehlt ein Teil der Wohnbevölkerung. Suche die Ursache in der Zuordnung der Tracts, schlage eine Korrektur vor und lass die Aufbereitung abbrechen, wenn zu wenig Bevölkerung zugeordnet ist.
      & s1\_\allowbreak{}daten.py, Code~3 \\
    \addlinespace
    Claude2026az & Anthropic Claude Opus 5 & 30.08.2026
      & Bevor ich Ergebnisse ansehe, will ich wissen, welche Stadtteile im Hold-out und welche in der Entwicklung liegen. Erstelle ein Profil beider Hälften mit Klassenverteilung und Zielgrößen, unabhängig von jedem Modellergebnis.
      & panelprofil.py \\
    \addlinespace
    Claude2026ba & Anthropic Claude Opus 5 & 02.09.2026
      & Probst et~al. zeigen, dass Verfahren unterschiedlich stark auf ihre Hyperparameter reagieren. Setze jeden getunten Parametersatz auf allen anderen Folds ein und fasse zusammen, wie weit die Güte dabei schwankt.
      & parameter\allowbreak{}sensitivitaet.py \\
    \addlinespace
    Claude2026bb & Anthropic Claude Opus 5 & 11.09.2026
      & Nach Altman und Bland ist ein nicht signifikanter Test kein Beleg für Gleichheit. Berechne für die gepaarten Vergleiche Effektstärke und Trennschärfe bei zehn Wiederholungsmitteln, mit derselben Formel für die Streuung wie in den Vergleichsskripten.
      & trennschaerfe.py \\
    \addlinespace
    Claude2026bc & Anthropic Claude Opus 5 & 14.09.2026
      & Setze die README für die Abgabe auf, gegliedert nach Umgebung, Daten, Ablauf, Reproduzierbarkeit und Aufbau. Den Ablauf gebe ich als Befehlsfolge in einer Datei vor. Alle weiteren Inhalte stammen aus Code und Konfiguration, dazu ein Hinweis, dass ein neuer Abruf die Rohdaten überschreibt und der Census-Schlüssel nicht beiliegt.
      & README.md \\
    \addlinespace
    Claude2026bd & Anthropic Claude Opus 5.5 & 24.09.2026
      & Gleiche die Kommentare aller Skripte mit dem Stand der Arbeit ab. Meine angehängte Liste nennt veraltete Begriffe, Zahlen und Kapitelverweise. Schlage je Datei die Änderungen vor, geändert werden nur Kommentare und Docstrings.
      & alle Skripte \\
    \addlinespace
    Claude2026be & Anthropic Claude Opus 5.5 & 24.09.2026
      & Bringe die Kommentare in allen Dateien in eine einheitliche Form nach meiner Vorlage: ein Kopf mit Zweck, Aufruf, Eingang und Ausgang, darunter kurze Stichpunkte zu den Entscheidungen. Interne Kürzel und Verweise auf Arbeitsnotizen entfallen.
      & alle Skripte \\
    \addlinespace
\end{xltabular}
% \tabquelle
\endgroup

% Benoetigt in der Praeambel: \usepackage{xltabular}
% Regeln: Projekt-Notiz KI-Verzeichnis_Regeln_2026-09-30
%   Spalten: Kennung, Modell, Datum, Prompt, Verwendung
%   Kennung: Claude2026a ... Claude2026z, Claude2026aa, Claude2026ab ...
%   Sortierung: nach Vorkommen, erst Gesamttext, dann Abschnitte
%               und Literatur, dann Code in der Reihenfolge der Entstehung
%               (Pipeline: Aufbereitung, Vorpruefung, Modelle, README,
%               Kommentare)
%   Prompts ohne Anfuehrungszeichen, ohne Kursivschrift
% ---------------------------------------------------
\newpage
\phantomsection
\section*{KI-Verzeichnis}
\addcontentsline{toc}{section}{KI-Verzeichnis}

\begingroup
\tabformat
\setlength{\tabcolsep}{4pt}
\begin{xltabular}{\textwidth}{@{}%
    >{\raggedright\arraybackslash\hyphenpenalty=10000}p{2.6cm}%
    >{\raggedright\arraybackslash}p{1.8cm}%
    >{\raggedright\arraybackslash\hyphenpenalty=10000}p{2.0cm}%
    >{\raggedright\arraybackslash}X%
    >{\raggedright\arraybackslash\hyphenpenalty=10000}p{2.8cm}@{}}
    % \caption*{Verzeichnis der verwendeten KI-Hilfsmittel}\\
    \toprule
    \textbf{Kennung} & \textbf{Modell} & \textbf{Datum} & \textbf{Prompt}
        & \textbf{Verwendung} \\
    \midrule
    \endfirsthead
    \toprule
    \textbf{Kennung} & \textbf{Modell} & \textbf{Datum} & \textbf{Prompt}
        & \textbf{Verwendung} \\
    \midrule
    \endhead
    \bottomrule
    \endlastfoot
    %
    Claude2026a & Anthropic Claude Opus 5.5 & 03.10.2026
      & Prüfe meine fertige Bachelorarbeit auf Fehler in Rechtschreibung, Zeichensetzung und Wortwahl und schlage für jede Stelle eine Korrektur vor. Inhalt, Aufbau und Argumentation bleiben unverändert. Ich entscheide bei jedem Vorschlag selbst, ob ich ihn übernehme.
      & gesamte Arbeit \\
    \addlinespace
    Claude2026b & Anthropic Claude Opus 5.5 & 03.10.2026
      & Schaue dir jeden Abschnitt an und gucke nach Schachtelsätzen, uneinheitlichen Fachbegriffen und schwerfälligen Formulierungen und schlage kürzere Alternativen vor. Halte dich dabei an meine beigefügten Stilregeln, etwa kurze Hauptsätze, ein Begriff je Sache und keine Verstärker.
      & gesamte Arbeit \\
    \addlinespace
    %
    Claude2026c & Anthropic Claude Opus 5 & 02.09.2026
      & Zu Kapitel 1 und 2 habe ich beim Gegenlesen eigene Anmerkungen notiert, sie hängen als Datei an. Ordne jede Anmerkung der betroffenen Textstelle zu und schlage je Stelle eine Änderung auf Satzebene vor.
      & Kap.~1, 2 \\
    \addlinespace
    Claude2026d & Anthropic Claude Opus 5 & 28.07.2026
      & Überarbeite den angehängten Absatz aus meinem Exposé für die Zielsetzung der Bachelorarbeit. Ziele und Gütemaße bleiben inhaltlich erhalten. Passe Begriffe und Formulierungen an meine angehängten Stichpunkte zum aktuellen Stand an.
      & Abschn.~1.2 \\
    \addlinespace
    Claude2026e & Anthropic Claude Opus 5 & 22.08.2026
      & In Shmuelis angehängtem Aufsatz geht es um den Unterschied zwischen erklärender und prognostischer Modellierung. Zieh mir die Definitionen beider Ansätze und die Stellen heraus, an denen der Aufsatz Folgen für die Modellwahl beschreibt, jeweils mit Seite.
      & Abschn.~2.2, 2.3, 2.7, 7.1 \\
    \addlinespace
    Claude2026f & Anthropic Claude Opus 5 & 28.08.2026
      & Suche mir frei zugängliche oder über die FOM erreichbare Standardwerke zu Grundlagen von Gütemaßen, Kreuzvalidierung und Hyperparametersuche. Nenne die Werke mit Kapitel oder Seite.
      & Abschn.~2.5 \\
    \addlinespace
    Claude2026g & Anthropic Claude Opus 5 & 02.09.2026
      & Lies im angehängten Text von Fahrmeir die Kapitel zu statistischen Tests und fasse zusammen, was dort zum Wilcoxon-Vorzeichen-Rang-Test, zur Gütefunktion und zum multiplen Testproblem steht. Ich brauche je Punkt die Seite.
      & Abschn.~2.5 \\
    \addlinespace
    Claude2026h & Anthropic Claude Opus 5 & 20.08.2026
      & Suche frei zugängliche Literatur zu Gütemaßen bei unausgewogenen Klassen, konkret zu Accuracy, Macro-F1 und mehrklassiger AUROC. Nenne je Werk die Seiten, auf denen das Thema behandelt wird, damit ich die Stellen nachlesen kann.
      & Abschn.~2.5, 6.2 \\
    \addlinespace
    Claude2026i & Anthropic Claude Opus 5 & 10.09.2026
      & Für jede der angehängten Studien brauche ich: Gegenstand, Datenbasis und zentralen Befund mit Seitenzahl. Formuliere den Befund in einem Satz.
      & Abschn.~2.7 \\
    \addlinespace
    Claude2026j & Anthropic Claude Opus 5 & 07.09.2026
      & Schlage mir den Rahmentext zur Quertabelle vor. Grundlage sind meine angehängten Stichpunkte zur Forschungslücke und zu den vier Punkten, in denen sich meine Arbeit von den bisherigen Studien unterscheidet. Der Text soll die Tabelle einordnen, nicht wiederholen.
      & Abschn.~2.7 \\
    \addlinespace
    Claude2026k & Anthropic Claude Opus 5 & 30.08.2026
      & Beschreibt das angehängte CRISP-DM-Handbuch von Chapman das Prozessmodell so, wie es der Aufsatz von Wirth und Hipp darstellt? Gleiche beide PDFs ab und nenne die Seiten, auf denen Chapman Phasen, Aufgaben und Abstraktionsebenen des Modells beschreibt.
      & Abschn.~3.1 \\
    \addlinespace
    Claude2026l & Anthropic Claude Opus 5 & 22.08.2026
      & Aus meinen angehängten Stichpunkten und der Datei der deskriptiven Auswertung soll ein Textvorschlag für die Datenbeschreibung und explorative Analyse entstehen. Übernimm nur Zahlen aus der Befunddatei, markiere jede Zahl, die du dort nicht findest.
      & Abschn.~4.2 \\
    \addlinespace
    Claude2026m & Anthropic Claude Opus 5 & 08.09.2026
      & Erkläre die beiden angehängten Codeausschnitte zur Verknüpfung der Zensusdaten mit den Stadtteilen und ordne ein, welche Rolle sie in der gesamten Datenaufbereitung spielen. Schlage mir dazu einen kurzen Text mit Bezug auf die Zeilennummern vor.
      & Abschn.~4.3 \\
    \addlinespace
    Claude2026n & Anthropic Claude Opus 5 & 09.09.2026
      & Fasse in Stichpunkten zusammen, was die angehängten Codeausschnitte zur Konstruktion der Merkmale tun und an welcher Stelle der Pipeline sie stehen. Nenne zu jedem Punkt die passenden Zeilennummern.
      & Abschn.~4.4 \\
    \addlinespace
    Claude2026o & Anthropic Claude Opus 5 & 09.09.2026
      & Welche Aufgabe haben die angehängten Codeausschnitte zu Fold-Zuteilung, Tests und Baselines im gesamten Ablauf der Auswertung? Mache mir einen kurzen Textvorschlag mit Bezug auf die Zeilennummern.
      & Abschn.~4.5 \\
    \addlinespace
    Claude2026p & Anthropic Claude Opus 5 & 25.08.2026
      & Im Anhang stehen die Verfahren und Tests, die für die Eignungsprüfung laufen: Ridge, Random Forest und XGBoost sowie RESET, Breusch-Pagan, Jarque-Bera und ein Überdispersionstest. Schreibe mir auf Basis der angehängten Arbeiten eine Beschreibung der Tests.
      & Abschn.~5.1 \\
    \addlinespace
    Claude2026q & Anthropic Claude Opus 5 & 07.09.2026
      & Schreibe mir zu jedem angehängten Codeausschnitt aus der Modellierung zwei bis drei Sätze: was der Code tut, wie er an die Datenaufbereitung anschließt und in welchen Zeilen das geschieht.
      & Abschn.~5.1, 5.2 \\
    \addlinespace
    Claude2026r & Anthropic Claude Opus 5 & 21.08.2026
      & Mache einen Vorschlag für eine Kürzung meines Ergebnisabschnitts um anderthalb Seiten, vor allem dort, wo der Text Tabellen nacherzählt oder Gesagtes wiederholt. Quellenverweise bleiben unangetastet, höchstens zwei Aussagen dürfen entfallen. Liste mir danach jede Streichung auf.
      & Abschn.~6.1, 6.2 \\
    \addlinespace
    Claude2026s & Anthropic Claude Opus 5 & 07.09.2026
      & Für den Abschnitt zum Aufwand fehlt mir die theoretische Laufzeitkomplexität der drei Verfahren. Schlage einen kurzen Textbaustein vor, der sich auf die angehängten Originalarbeiten stützt, und markiere jede Seitenangabe, die ich noch am Volltext prüfen muss.
      & Abschn.~6.4 \\
    \addlinespace
    Claude2026t & Anthropic Claude Opus 5 & 26.08.2026
      & Hier sind meine Stichpunkte zur Diskussion als TXT-Datei, dazu die Ergebnistabellen aus Kapitel~6. Schlage mir für Einordnung, Implikationen und Limitationen je einen Textvorschlag vor, der meine Gliederung übernimmt und nichts aus Kapitel~6 wiederholt.
      & Abschn.~7.1--7.3 \\
    \addlinespace
    Claude2026u & Anthropic Claude Opus 5 & 26.08.2026
      & Der angehängte Diskussionsentwurf unterschätzt aus meiner Sicht, was die Merkmale leisten. Prüfe Einordnung und Implikationen gegen meine Ergebnistabelle, ergänze einen Vergleich mit einer weiteren Referenz aus dieser Tabelle und schlage überarbeitete Absätze vor.
      & Abschn.~7.1, 7.2 \\
    \addlinespace
    Claude2026v & Anthropic Claude Opus 5 & 27.07.2026
      & Lege nach meinen angehängten Festlegungen eine zentrale Konfiguration für die Aufbereitung an: Analysezeitraum, Plausibilitätsfenster der Antwortzeit, ausgeschlossene Stadtteile und Schalter für die Downloads. Alle Skripte sollen diese Werte nur lesen und nie selbst setzen.
      & config.py \\
    \addlinespace
    Claude2026w & Anthropic Claude Opus 5 & 27.07.2026
      & Schreibe den Abruf der Rohdaten über die Schnittstellen von DataSF und dem Census Bureau. Welche Datensätze, Felder und Filter gebraucht werden, steht in meiner Quellenliste als Datei. Jede Quelle wird unverändert unter data/raw abgelegt.
      & s1\_\allowbreak{}daten.py \\
    \addlinespace
    Claude2026x & Anthropic Claude Opus 5 & 28.07.2026
      & Die Sozialdaten liegen je Census Tract vor. Ordne sie über die angehängte Zuordnungstabelle den Stadtteilen zu und gib je Jahrgang die Trefferquote aus.
      & s1\_\allowbreak{}daten.py, Code~3 \\
    \addlinespace
    Claude2026y & Anthropic Claude Opus 5 & 28.07.2026
      & Kaufman et~al. beschreiben Leakage auch durch Werte, die zum Vorhersagezeitpunkt noch gar nicht bekannt waren. Ordne deshalb jedem Monat nur den ACS-Jahrgang zu, der damals schon veröffentlicht war. Den Versatz je Jahrgang habe ich angehängt.
      & s1\_\allowbreak{}daten.py, Code~4 \\
    \addlinespace
    Claude2026z & Anthropic Claude Opus 5 & 28.07.2026
      & Bilde die Anteilswerte aus meiner angehängten Merkmalsliste. Ist ein Nenner null, soll ein fehlender Wert entstehen.
      & s1\_\allowbreak{}daten.py, Code~6 \\
    \addlinespace
    Claude2026aa & Anthropic Claude Opus 5 & 29.07.2026
      & Für die Kriminalität habe ich mich nach Brantingham und Brantingham für einen Standortquotienten entschieden. Implementiere ihn als relativen Index je Stadtteil und Monat, bezogen auf die gesamte Stadt im selben Monat.
      & s1\_\allowbreak{}daten.py, Code~7 \\
    \addlinespace
    Claude2026ab & Anthropic Claude Opus 5 & 29.07.2026
      & Aggregiere die Einsätze auf Stadtteil und Monat, einmal als Anzahl und einmal als überwiegende Einsatzart. Jede Kombination soll als Zeile vorhanden sein, jährliche Werte werden nach meiner angehängten Regel vorwärts gefüllt.
      & s2\_\allowbreak{}datensaetze.py, Code~5 \\
    \addlinespace
    Claude2026ac & Anthropic Claude Opus 5 & 29.07.2026
      & Kodiere den Kalendermonat als Sinus und Kosinus und ergänze die verzögerten Werte aus meiner Merkmalsliste, nur aus abgeschlossenen Vormonaten.
      & s2\_\allowbreak{}datensaetze.py, Code~8 \\
    \addlinespace
    Claude2026ad & Anthropic Claude Opus 5 & 30.07.2026
      & Roberts et~al. empfehlen bei räumlich abhängigen Daten eine gruppierte Kreuzvalidierung. Teile deshalb ganze Stadtteile auf fünf Folds und ein Hold-out auf, geschichtet nach brand-dominierten Monaten und Bevölkerung. Die Zuteilung steht als Spalte im Datensatz.
      & s2\_\allowbreak{}datensaetze.py, Code~10 \\
    \addlinespace
    Claude2026ae & Anthropic Claude Opus 5 & 30.07.2026
      & Baue eine Prüfung ein, die nach dem Bau des Panels abbricht, wenn Stadtteile, Zeitraum oder Wohnbevölkerung nicht zu meinen angehängten Sollwerten passen. Die Fehlermeldung soll den betroffenen Stadtteil und den gemessenen Wert nennen.
      & s2\_\allowbreak{}datensaetze.py, Code~1 \\
    \addlinespace
    Claude2026af & Anthropic Claude Opus 5 & 31.07.2026
      & Wie mit meinem Betreuer abgestimmt, messen sich die Verfahren an zwei Baselines. Implementiere eine Referenz ohne Merkmale und ein Poisson-GLM mit der Bevölkerung als Offset nach Winkelmann, für die Einsatzart ein multinomiales Logit.
      & v1\_\allowbreak{}baselines.py, Code~12, 13 \\
    \addlinespace
    Claude2026ag & Anthropic Claude Opus 5 & 03.08.2026
      & Setze die angehängten Eignungsprüfungen, die ich vor der Modellierung festgelegt habe, als ein Skript mit einem Bericht am Ende um. Dazu gehören die Überdispersion als Hilfsregression, die Linearität grafisch vor dem Einsatz von Ridge, der RESET nach Ramsey und je Verfahren eine Tabelle mit Anforderung, Teststatistik und p-Wert.
      & v2\_\allowbreak{}eignung.py, Code~14 \\
    \addlinespace
    Claude2026ah & Anthropic Claude Opus 5 & 03.08.2026
      & Fasse die Aufbereitung in einem einzigen Befehl zusammen, der beide Schritte nacheinander ausführt, bei einem Fehler abbricht und am Ende Zeilen und Spalten beider Dateien ausgibt. Dasselbe brauche ich für die Vorprüfung: erst die Baselines, danach die Eignungsprüfung, die deren Werte nur liest.
      & build.py, run.py \\
    \addlinespace
    Claude2026ai & Anthropic Claude Opus 5 & 04.08.2026
      & Lege das Skript für den Mengenvergleich in Phasen an, die ich nacheinander fülle: Tuning, Bewertung, Aggregation und Vergleich. Ridge rechnet nach meiner Methodik auf $\log(1+y)$, Random Forest und XGBoost mit zählgerechter Verlustfunktion. Ein einzelner Lauf passt auf der Rate an, rechnet auf Einsätze zurück und misst die Zeit.
      & m02\_\allowbreak{}menge.py, Code~15, 16 \\
    \addlinespace
    Claude2026aj & Anthropic Claude Opus 5 & 05.08.2026
      & Übertrage den Aufbau des Mengenvergleichs auf die Einsatzart mit Random Forest und XGBoost. Die vier Klassen sollen global einheitlich kodiert und ausgeglichen gewichtet werden, Hauptmaß ist nach Opitz und Burst der Macro-F1.
      & m03\_\allowbreak{}struktur.py, Code~17 \\
    \addlinespace
    Claude2026ak & Anthropic Claude Opus 5 & 06.08.2026
      & Die zehn Wiederholungen verschieben bisher nur die Nummern der Folds, die Zusammensetzung bleibt gleich, und das Hold-out wandert mit. Korrigiere das so, dass innerhalb der Rangblöcke meiner Schichtung neu gemischt wird und das Hold-out fest bleibt.
      & v0\_\allowbreak{}aufteilung.py \\
    \addlinespace
    Claude2026al & Anthropic Claude Opus 5 & 06.08.2026
      & Nach Bergstra und Bengio nehme ich eine Zufallssuche. Setze sie als innere, nach Stadtteilen gruppierte Kreuzvalidierung mit meinem angehängten Suchraum um.
      & m02\_\allowbreak{}menge.py, Code~18 \\
    \addlinespace
    Claude2026am & Anthropic Claude Opus 5 & 07.08.2026
      & Hurlbert beschreibt Pseudoreplikation, wenn abhängige Messungen als unabhängig gezählt werden. Mittle die Gütemaße deshalb zweistufig, erst über die Folds je Wiederholung und dann über die zehn Wiederholungen, und gib beide Streuungen aus.
      & m02\_\allowbreak{}menge.py, Code~11, 19 \\
    \addlinespace
    Claude2026an & Anthropic Claude Opus 5 & 07.08.2026
      & Die Eignungsprüfung verwirft die lineare Form nur innerhalb der Stichprobe. Prüfe in einem eigenen Skript, ob Quadrat- und Interaktionsterme im Poisson-GLM auch auf unbekannten Stadtteilen besser vorhersagen, mit denselben Folds wie im Vergleich.
      & v3\_\allowbreak{}spezifikation.py \\
    \addlinespace
    Claude2026ao & Anthropic Claude Opus 5 & 08.08.2026
      & Vergleiche jedes Verfahren gepaart mit dem Wilcoxon-Test gegen die Stufe-2-Baseline und korrigiere nach Holm. Die Schlussbewertung auf dem Hold-out läuft nur mit ausdrücklichem Argument, einmalig und mit den Parametern aus dem Tuning.
      & m02\_\allowbreak{}menge.py, m03\_\allowbreak{}struktur.py \\
    \addlinespace
    Claude2026ap & Anthropic Claude Opus 5 & 08.08.2026
      & Wiederholung null dient zugleich dem Tuning und der Bewertung. Beziffere, wie viel Vorsprung die Verfahren dadurch bekommen.
      & m02\_\allowbreak{}menge.py, m03\_\allowbreak{}struktur.py \\
    \addlinespace
    Claude2026aq & Anthropic Claude Opus 5 & 11.08.2026
      & Ich möchte die fertigen Dateien prüfen und nicht den Code. Schreibe dazu eine eigene Testdatei nach meiner angehängten Liste, gegliedert nach Panel, Aufteilung, Merkmalen und Struktur. Darunter fallen keine Ergebnisvariable als Merkmal im Sinne von Kaufman et~al. und kein Stadtteil zugleich in Training und Test. Sie soll auch ohne pytest laufen.
      & test\_\allowbreak{}aufbereitung.py, Code~2, 9 \\
    \addlinespace
    Claude2026ar & Anthropic Claude Opus 5 & 13.08.2026
      & SHAP-Werte nach Lundberg und Lee zeigen nur, wie ein Modell seine Aufmerksamkeit verteilt. Ergänze deshalb eine Ablation je Faktorgruppe aus meiner Gruppierung, bei der das Modell ohne die Gruppe neu angepasst und bewertet wird.
      & m04\_\allowbreak{}shap.py \\
    \addlinespace
    Claude2026as & Anthropic Claude Opus 5 & 14.08.2026
      & Sorge dafür, dass ein zweiter Lauf dieselben Zahlen liefert: ein fester Startwert für alle Zufallsschritte, Bewertung auf einem Kern und eine Datei mit den genauen Paketversionen. Die Konstanten der Modellierung kommen in eine eigene Konfiguration.
      & config\_\allowbreak{}modelle.py, requirements\_\allowbreak{}lauf.txt \\
    \addlinespace
    Claude2026at & Anthropic Claude Opus 5 & 15.08.2026
      & Ob das Budget der Hyperparametersuche reicht, will ich wie von Bischl et~al. empfohlen prüfen. Wiederhole die Suche mit steigender Zahl an Versuchen, zeige den Verlauf als Kurve und ergänze einen kurzen Probelauf.
      & suchdiagnose.py \\
    \addlinespace
    Claude2026au & Anthropic Claude Opus 5 & 17.08.2026
      & Ein Macro-F1 lässt sich ohne Obergrenze schlecht einordnen. Schätze mit meinem angehängten Ansatz, wie gut die Einsatzart mit diesen Merkmalen höchstens vorhersagbar ist, und setze die Werte aus dem Vergleich ins Verhältnis dazu.
      & v4\_\allowbreak{}decke.py \\
    \addlinespace
    Claude2026av & Anthropic Claude Opus 5 & 20.08.2026
      & Zum Data Understanding nach Chapman et~al. brauche ich zwei getrennte Skripte: eine Merkmalstabelle mit Skalenniveau, Einheit und Quelle aus meiner Vorlage und eine deskriptive Auswertung mit Verteilung, Varianzzerlegung und Korrelationen.
      & codebook.py, deskriptiv.py \\
    \addlinespace
    Claude2026aw & Anthropic Claude Opus 5 & 23.08.2026
      & Der Census-Schlüssel steht im Klartext in der Konfiguration. Lies ihn stattdessen aus einer Umgebungsvariable.
      & config.py \\
    \addlinespace
    Claude2026ax & Anthropic Claude Opus 5 & 25.08.2026
      & Stelle aus den Rohquellen die Befunde zusammen, aus denen ein Eingriff in der Aufbereitung folgt, etwa Dubletten, fehlende Baujahre oder Lücken in der Abdeckung. Meine Liste der Eingriffe liegt bei, alles andere bleibt draußen.
      & rohbefunde.py \\
    \addlinespace
    Claude2026ay & Anthropic Claude Opus 5 & 29.08.2026
      & In den ACS-Jahrgängen mit älteren Zensusgrenzen fehlt ein Teil der Wohnbevölkerung. Suche die Ursache in der Zuordnung der Tracts, schlage eine Korrektur vor und lass die Aufbereitung abbrechen, wenn zu wenig Bevölkerung zugeordnet ist.
      & s1\_\allowbreak{}daten.py, Code~3 \\
    \addlinespace
    Claude2026az & Anthropic Claude Opus 5 & 30.08.2026
      & Bevor ich Ergebnisse ansehe, will ich wissen, welche Stadtteile im Hold-out und welche in der Entwicklung liegen. Erstelle ein Profil beider Hälften mit Klassenverteilung und Zielgrößen, unabhängig von jedem Modellergebnis.
      & panelprofil.py \\
    \addlinespace
    Claude2026ba & Anthropic Claude Opus 5 & 02.09.2026
      & Probst et~al. zeigen, dass Verfahren unterschiedlich stark auf ihre Hyperparameter reagieren. Setze jeden getunten Parametersatz auf allen anderen Folds ein und fasse zusammen, wie weit die Güte dabei schwankt.
      & parameter\allowbreak{}sensitivitaet.py \\
    \addlinespace
    Claude2026bb & Anthropic Claude Opus 5 & 11.09.2026
      & Nach Altman und Bland ist ein nicht signifikanter Test kein Beleg für Gleichheit. Berechne für die gepaarten Vergleiche Effektstärke und Trennschärfe bei zehn Wiederholungsmitteln, mit derselben Formel für die Streuung wie in den Vergleichsskripten.
      & trennschaerfe.py \\
    \addlinespace
    Claude2026bc & Anthropic Claude Opus 5 & 14.09.2026
      & Setze die README für die Abgabe auf, gegliedert nach Umgebung, Daten, Ablauf, Reproduzierbarkeit und Aufbau. Den Ablauf gebe ich als Befehlsfolge in einer Datei vor. Alle weiteren Inhalte stammen aus Code und Konfiguration, dazu ein Hinweis, dass ein neuer Abruf die Rohdaten überschreibt und der Census-Schlüssel nicht beiliegt.
      & README.md \\
    \addlinespace
    Claude2026bd & Anthropic Claude Opus 5.5 & 24.09.2026
      & Gleiche die Kommentare aller Skripte mit dem Stand der Arbeit ab. Meine angehängte Liste nennt veraltete Begriffe, Zahlen und Kapitelverweise. Schlage je Datei die Änderungen vor, geändert werden nur Kommentare und Docstrings.
      & alle Skripte \\
    \addlinespace
    Claude2026be & Anthropic Claude Opus 5.5 & 24.09.2026
      & Bringe die Kommentare in allen Dateien in eine einheitliche Form nach meiner Vorlage: ein Kopf mit Zweck, Aufruf, Eingang und Ausgang, darunter kurze Stichpunkte zu den Entscheidungen. Interne Kürzel und Verweise auf Arbeitsnotizen entfallen.
      & alle Skripte \\
    \addlinespace
\end{xltabular}
% \tabquelle
\endgroup

% Benoetigt in der Praeambel: \usepackage{xltabular}
% Regeln: Projekt-Notiz KI-Verzeichnis_Regeln_2026-09-30
%   Spalten: Kennung, Modell, Datum, Prompt, Verwendung
%   Kennung: Claude2026a ... Claude2026z, Claude2026aa, Claude2026ab ...
%   Sortierung: nach Vorkommen, erst Gesamttext, dann Abschnitte
%               und Literatur, dann Code in der Reihenfolge der Entstehung
%               (Pipeline: Aufbereitung, Vorpruefung, Modelle, README,
%               Kommentare)
%   Prompts ohne Anfuehrungszeichen, ohne Kursivschrift
% ---------------------------------------------------
\newpage
\phantomsection
\section*{KI-Verzeichnis}
\addcontentsline{toc}{section}{KI-Verzeichnis}

\begingroup
\tabformat
\setlength{\tabcolsep}{4pt}
\begin{xltabular}{\textwidth}{@{}%
    >{\raggedright\arraybackslash\hyphenpenalty=10000}p{2.6cm}%
    >{\raggedright\arraybackslash}p{1.8cm}%
    >{\raggedright\arraybackslash\hyphenpenalty=10000}p{2.0cm}%
    >{\raggedright\arraybackslash}X%
    >{\raggedright\arraybackslash\hyphenpenalty=10000}p{2.8cm}@{}}
    % \caption*{Verzeichnis der verwendeten KI-Hilfsmittel}\\
    \toprule
    \textbf{Kennung} & \textbf{Modell} & \textbf{Datum} & \textbf{Prompt}
        & \textbf{Verwendung} \\
    \midrule
    \endfirsthead
    \toprule
    \textbf{Kennung} & \textbf{Modell} & \textbf{Datum} & \textbf{Prompt}
        & \textbf{Verwendung} \\
    \midrule
    \endhead
    \bottomrule
    \endlastfoot
    %
    Claude2026a & Anthropic Claude Opus 5.5 & 03.10.2026
      & Prüfe meine fertige Bachelorarbeit auf Fehler in Rechtschreibung, Zeichensetzung und Wortwahl und schlage für jede Stelle eine Korrektur vor. Inhalt, Aufbau und Argumentation bleiben unverändert. Ich entscheide bei jedem Vorschlag selbst, ob ich ihn übernehme.
      & gesamte Arbeit \\
    \addlinespace
    Claude2026b & Anthropic Claude Opus 5.5 & 03.10.2026
      & Schaue dir jeden Abschnitt an und gucke nach Schachtelsätzen, uneinheitlichen Fachbegriffen und schwerfälligen Formulierungen und schlage kürzere Alternativen vor. Halte dich dabei an meine beigefügten Stilregeln, etwa kurze Hauptsätze, ein Begriff je Sache und keine Verstärker.
      & gesamte Arbeit \\
    \addlinespace
    %
    Claude2026c & Anthropic Claude Opus 5 & 02.09.2026
      & Zu Kapitel 1 und 2 habe ich beim Gegenlesen eigene Anmerkungen notiert, sie hängen als Datei an. Ordne jede Anmerkung der betroffenen Textstelle zu und schlage je Stelle eine Änderung auf Satzebene vor.
      & Kap.~1, 2 \\
    \addlinespace
    Claude2026d & Anthropic Claude Opus 5 & 28.07.2026
      & Überarbeite den angehängten Absatz aus meinem Exposé für die Zielsetzung der Bachelorarbeit. Ziele und Gütemaße bleiben inhaltlich erhalten. Passe Begriffe und Formulierungen an meine angehängten Stichpunkte zum aktuellen Stand an.
      & Abschn.~1.2 \\
    \addlinespace
    Claude2026e & Anthropic Claude Opus 5 & 22.08.2026
      & In Shmuelis angehängtem Aufsatz geht es um den Unterschied zwischen erklärender und prognostischer Modellierung. Zieh mir die Definitionen beider Ansätze und die Stellen heraus, an denen der Aufsatz Folgen für die Modellwahl beschreibt, jeweils mit Seite.
      & Abschn.~2.2, 2.3, 2.7, 7.1 \\
    \addlinespace
    Claude2026f & Anthropic Claude Opus 5 & 28.08.2026
      & Suche mir frei zugängliche oder über die FOM erreichbare Standardwerke zu Grundlagen von Gütemaßen, Kreuzvalidierung und Hyperparametersuche. Nenne die Werke mit Kapitel oder Seite.
      & Abschn.~2.5 \\
    \addlinespace
    Claude2026g & Anthropic Claude Opus 5 & 02.09.2026
      & Lies im angehängten Text von Fahrmeir die Kapitel zu statistischen Tests und fasse zusammen, was dort zum Wilcoxon-Vorzeichen-Rang-Test, zur Gütefunktion und zum multiplen Testproblem steht. Ich brauche je Punkt die Seite.
      & Abschn.~2.5 \\
    \addlinespace
    Claude2026h & Anthropic Claude Opus 5 & 20.08.2026
      & Suche frei zugängliche Literatur zu Gütemaßen bei unausgewogenen Klassen, konkret zu Accuracy, Macro-F1 und mehrklassiger AUROC. Nenne je Werk die Seiten, auf denen das Thema behandelt wird, damit ich die Stellen nachlesen kann.
      & Abschn.~2.5, 6.2 \\
    \addlinespace
    Claude2026i & Anthropic Claude Opus 5 & 10.09.2026
      & Für jede der angehängten Studien brauche ich: Gegenstand, Datenbasis und zentralen Befund mit Seitenzahl. Formuliere den Befund in einem Satz.
      & Abschn.~2.7 \\
    \addlinespace
    Claude2026j & Anthropic Claude Opus 5 & 07.09.2026
      & Schlage mir den Rahmentext zur Quertabelle vor. Grundlage sind meine angehängten Stichpunkte zur Forschungslücke und zu den vier Punkten, in denen sich meine Arbeit von den bisherigen Studien unterscheidet. Der Text soll die Tabelle einordnen, nicht wiederholen.
      & Abschn.~2.7 \\
    \addlinespace
    Claude2026k & Anthropic Claude Opus 5 & 30.08.2026
      & Beschreibt das angehängte CRISP-DM-Handbuch von Chapman das Prozessmodell so, wie es der Aufsatz von Wirth und Hipp darstellt? Gleiche beide PDFs ab und nenne die Seiten, auf denen Chapman Phasen, Aufgaben und Abstraktionsebenen des Modells beschreibt.
      & Abschn.~3.1 \\
    \addlinespace
    Claude2026l & Anthropic Claude Opus 5 & 22.08.2026
      & Aus meinen angehängten Stichpunkten und der Datei der deskriptiven Auswertung soll ein Textvorschlag für die Datenbeschreibung und explorative Analyse entstehen. Übernimm nur Zahlen aus der Befunddatei, markiere jede Zahl, die du dort nicht findest.
      & Abschn.~4.2 \\
    \addlinespace
    Claude2026m & Anthropic Claude Opus 5 & 08.09.2026
      & Erkläre die beiden angehängten Codeausschnitte zur Verknüpfung der Zensusdaten mit den Stadtteilen und ordne ein, welche Rolle sie in der gesamten Datenaufbereitung spielen. Schlage mir dazu einen kurzen Text mit Bezug auf die Zeilennummern vor.
      & Abschn.~4.3 \\
    \addlinespace
    Claude2026n & Anthropic Claude Opus 5 & 09.09.2026
      & Fasse in Stichpunkten zusammen, was die angehängten Codeausschnitte zur Konstruktion der Merkmale tun und an welcher Stelle der Pipeline sie stehen. Nenne zu jedem Punkt die passenden Zeilennummern.
      & Abschn.~4.4 \\
    \addlinespace
    Claude2026o & Anthropic Claude Opus 5 & 09.09.2026
      & Welche Aufgabe haben die angehängten Codeausschnitte zu Fold-Zuteilung, Tests und Baselines im gesamten Ablauf der Auswertung? Mache mir einen kurzen Textvorschlag mit Bezug auf die Zeilennummern.
      & Abschn.~4.5 \\
    \addlinespace
    Claude2026p & Anthropic Claude Opus 5 & 25.08.2026
      & Im Anhang stehen die Verfahren und Tests, die für die Eignungsprüfung laufen: Ridge, Random Forest und XGBoost sowie RESET, Breusch-Pagan, Jarque-Bera und ein Überdispersionstest. Schreibe mir auf Basis der angehängten Arbeiten eine Beschreibung der Tests.
      & Abschn.~5.1 \\
    \addlinespace
    Claude2026q & Anthropic Claude Opus 5 & 07.09.2026
      & Schreibe mir zu jedem angehängten Codeausschnitt aus der Modellierung zwei bis drei Sätze: was der Code tut, wie er an die Datenaufbereitung anschließt und in welchen Zeilen das geschieht.
      & Abschn.~5.1, 5.2 \\
    \addlinespace
    Claude2026r & Anthropic Claude Opus 5 & 21.08.2026
      & Mache einen Vorschlag für eine Kürzung meines Ergebnisabschnitts um anderthalb Seiten, vor allem dort, wo der Text Tabellen nacherzählt oder Gesagtes wiederholt. Quellenverweise bleiben unangetastet, höchstens zwei Aussagen dürfen entfallen. Liste mir danach jede Streichung auf.
      & Abschn.~6.1, 6.2 \\
    \addlinespace
    Claude2026s & Anthropic Claude Opus 5 & 07.09.2026
      & Für den Abschnitt zum Aufwand fehlt mir die theoretische Laufzeitkomplexität der drei Verfahren. Schlage einen kurzen Textbaustein vor, der sich auf die angehängten Originalarbeiten stützt, und markiere jede Seitenangabe, die ich noch am Volltext prüfen muss.
      & Abschn.~6.4 \\
    \addlinespace
    Claude2026t & Anthropic Claude Opus 5 & 26.08.2026
      & Hier sind meine Stichpunkte zur Diskussion als TXT-Datei, dazu die Ergebnistabellen aus Kapitel~6. Schlage mir für Einordnung, Implikationen und Limitationen je einen Textvorschlag vor, der meine Gliederung übernimmt und nichts aus Kapitel~6 wiederholt.
      & Abschn.~7.1--7.3 \\
    \addlinespace
    Claude2026u & Anthropic Claude Opus 5 & 26.08.2026
      & Der angehängte Diskussionsentwurf unterschätzt aus meiner Sicht, was die Merkmale leisten. Prüfe Einordnung und Implikationen gegen meine Ergebnistabelle, ergänze einen Vergleich mit einer weiteren Referenz aus dieser Tabelle und schlage überarbeitete Absätze vor.
      & Abschn.~7.1, 7.2 \\
    \addlinespace
    Claude2026v & Anthropic Claude Opus 5 & 27.07.2026
      & Lege nach meinen angehängten Festlegungen eine zentrale Konfiguration für die Aufbereitung an: Analysezeitraum, Plausibilitätsfenster der Antwortzeit, ausgeschlossene Stadtteile und Schalter für die Downloads. Alle Skripte sollen diese Werte nur lesen und nie selbst setzen.
      & config.py \\
    \addlinespace
    Claude2026w & Anthropic Claude Opus 5 & 27.07.2026
      & Schreibe den Abruf der Rohdaten über die Schnittstellen von DataSF und dem Census Bureau. Welche Datensätze, Felder und Filter gebraucht werden, steht in meiner Quellenliste als Datei. Jede Quelle wird unverändert unter data/raw abgelegt.
      & s1\_\allowbreak{}daten.py \\
    \addlinespace
    Claude2026x & Anthropic Claude Opus 5 & 28.07.2026
      & Die Sozialdaten liegen je Census Tract vor. Ordne sie über die angehängte Zuordnungstabelle den Stadtteilen zu und gib je Jahrgang die Trefferquote aus.
      & s1\_\allowbreak{}daten.py, Code~3 \\
    \addlinespace
    Claude2026y & Anthropic Claude Opus 5 & 28.07.2026
      & Kaufman et~al. beschreiben Leakage auch durch Werte, die zum Vorhersagezeitpunkt noch gar nicht bekannt waren. Ordne deshalb jedem Monat nur den ACS-Jahrgang zu, der damals schon veröffentlicht war. Den Versatz je Jahrgang habe ich angehängt.
      & s1\_\allowbreak{}daten.py, Code~4 \\
    \addlinespace
    Claude2026z & Anthropic Claude Opus 5 & 28.07.2026
      & Bilde die Anteilswerte aus meiner angehängten Merkmalsliste. Ist ein Nenner null, soll ein fehlender Wert entstehen.
      & s1\_\allowbreak{}daten.py, Code~6 \\
    \addlinespace
    Claude2026aa & Anthropic Claude Opus 5 & 29.07.2026
      & Für die Kriminalität habe ich mich nach Brantingham und Brantingham für einen Standortquotienten entschieden. Implementiere ihn als relativen Index je Stadtteil und Monat, bezogen auf die gesamte Stadt im selben Monat.
      & s1\_\allowbreak{}daten.py, Code~7 \\
    \addlinespace
    Claude2026ab & Anthropic Claude Opus 5 & 29.07.2026
      & Aggregiere die Einsätze auf Stadtteil und Monat, einmal als Anzahl und einmal als überwiegende Einsatzart. Jede Kombination soll als Zeile vorhanden sein, jährliche Werte werden nach meiner angehängten Regel vorwärts gefüllt.
      & s2\_\allowbreak{}datensaetze.py, Code~5 \\
    \addlinespace
    Claude2026ac & Anthropic Claude Opus 5 & 29.07.2026
      & Kodiere den Kalendermonat als Sinus und Kosinus und ergänze die verzögerten Werte aus meiner Merkmalsliste, nur aus abgeschlossenen Vormonaten.
      & s2\_\allowbreak{}datensaetze.py, Code~8 \\
    \addlinespace
    Claude2026ad & Anthropic Claude Opus 5 & 30.07.2026
      & Roberts et~al. empfehlen bei räumlich abhängigen Daten eine gruppierte Kreuzvalidierung. Teile deshalb ganze Stadtteile auf fünf Folds und ein Hold-out auf, geschichtet nach brand-dominierten Monaten und Bevölkerung. Die Zuteilung steht als Spalte im Datensatz.
      & s2\_\allowbreak{}datensaetze.py, Code~10 \\
    \addlinespace
    Claude2026ae & Anthropic Claude Opus 5 & 30.07.2026
      & Baue eine Prüfung ein, die nach dem Bau des Panels abbricht, wenn Stadtteile, Zeitraum oder Wohnbevölkerung nicht zu meinen angehängten Sollwerten passen. Die Fehlermeldung soll den betroffenen Stadtteil und den gemessenen Wert nennen.
      & s2\_\allowbreak{}datensaetze.py, Code~1 \\
    \addlinespace
    Claude2026af & Anthropic Claude Opus 5 & 31.07.2026
      & Wie mit meinem Betreuer abgestimmt, messen sich die Verfahren an zwei Baselines. Implementiere eine Referenz ohne Merkmale und ein Poisson-GLM mit der Bevölkerung als Offset nach Winkelmann, für die Einsatzart ein multinomiales Logit.
      & v1\_\allowbreak{}baselines.py, Code~12, 13 \\
    \addlinespace
    Claude2026ag & Anthropic Claude Opus 5 & 03.08.2026
      & Setze die angehängten Eignungsprüfungen, die ich vor der Modellierung festgelegt habe, als ein Skript mit einem Bericht am Ende um. Dazu gehören die Überdispersion als Hilfsregression, die Linearität grafisch vor dem Einsatz von Ridge, der RESET nach Ramsey und je Verfahren eine Tabelle mit Anforderung, Teststatistik und p-Wert.
      & v2\_\allowbreak{}eignung.py, Code~14 \\
    \addlinespace
    Claude2026ah & Anthropic Claude Opus 5 & 03.08.2026
      & Fasse die Aufbereitung in einem einzigen Befehl zusammen, der beide Schritte nacheinander ausführt, bei einem Fehler abbricht und am Ende Zeilen und Spalten beider Dateien ausgibt. Dasselbe brauche ich für die Vorprüfung: erst die Baselines, danach die Eignungsprüfung, die deren Werte nur liest.
      & build.py, run.py \\
    \addlinespace
    Claude2026ai & Anthropic Claude Opus 5 & 04.08.2026
      & Lege das Skript für den Mengenvergleich in Phasen an, die ich nacheinander fülle: Tuning, Bewertung, Aggregation und Vergleich. Ridge rechnet nach meiner Methodik auf $\log(1+y)$, Random Forest und XGBoost mit zählgerechter Verlustfunktion. Ein einzelner Lauf passt auf der Rate an, rechnet auf Einsätze zurück und misst die Zeit.
      & m02\_\allowbreak{}menge.py, Code~15, 16 \\
    \addlinespace
    Claude2026aj & Anthropic Claude Opus 5 & 05.08.2026
      & Übertrage den Aufbau des Mengenvergleichs auf die Einsatzart mit Random Forest und XGBoost. Die vier Klassen sollen global einheitlich kodiert und ausgeglichen gewichtet werden, Hauptmaß ist nach Opitz und Burst der Macro-F1.
      & m03\_\allowbreak{}struktur.py, Code~17 \\
    \addlinespace
    Claude2026ak & Anthropic Claude Opus 5 & 06.08.2026
      & Die zehn Wiederholungen verschieben bisher nur die Nummern der Folds, die Zusammensetzung bleibt gleich, und das Hold-out wandert mit. Korrigiere das so, dass innerhalb der Rangblöcke meiner Schichtung neu gemischt wird und das Hold-out fest bleibt.
      & v0\_\allowbreak{}aufteilung.py \\
    \addlinespace
    Claude2026al & Anthropic Claude Opus 5 & 06.08.2026
      & Nach Bergstra und Bengio nehme ich eine Zufallssuche. Setze sie als innere, nach Stadtteilen gruppierte Kreuzvalidierung mit meinem angehängten Suchraum um.
      & m02\_\allowbreak{}menge.py, Code~18 \\
    \addlinespace
    Claude2026am & Anthropic Claude Opus 5 & 07.08.2026
      & Hurlbert beschreibt Pseudoreplikation, wenn abhängige Messungen als unabhängig gezählt werden. Mittle die Gütemaße deshalb zweistufig, erst über die Folds je Wiederholung und dann über die zehn Wiederholungen, und gib beide Streuungen aus.
      & m02\_\allowbreak{}menge.py, Code~11, 19 \\
    \addlinespace
    Claude2026an & Anthropic Claude Opus 5 & 07.08.2026
      & Die Eignungsprüfung verwirft die lineare Form nur innerhalb der Stichprobe. Prüfe in einem eigenen Skript, ob Quadrat- und Interaktionsterme im Poisson-GLM auch auf unbekannten Stadtteilen besser vorhersagen, mit denselben Folds wie im Vergleich.
      & v3\_\allowbreak{}spezifikation.py \\
    \addlinespace
    Claude2026ao & Anthropic Claude Opus 5 & 08.08.2026
      & Vergleiche jedes Verfahren gepaart mit dem Wilcoxon-Test gegen die Stufe-2-Baseline und korrigiere nach Holm. Die Schlussbewertung auf dem Hold-out läuft nur mit ausdrücklichem Argument, einmalig und mit den Parametern aus dem Tuning.
      & m02\_\allowbreak{}menge.py, m03\_\allowbreak{}struktur.py \\
    \addlinespace
    Claude2026ap & Anthropic Claude Opus 5 & 08.08.2026
      & Wiederholung null dient zugleich dem Tuning und der Bewertung. Beziffere, wie viel Vorsprung die Verfahren dadurch bekommen.
      & m02\_\allowbreak{}menge.py, m03\_\allowbreak{}struktur.py \\
    \addlinespace
    Claude2026aq & Anthropic Claude Opus 5 & 11.08.2026
      & Ich möchte die fertigen Dateien prüfen und nicht den Code. Schreibe dazu eine eigene Testdatei nach meiner angehängten Liste, gegliedert nach Panel, Aufteilung, Merkmalen und Struktur. Darunter fallen keine Ergebnisvariable als Merkmal im Sinne von Kaufman et~al. und kein Stadtteil zugleich in Training und Test. Sie soll auch ohne pytest laufen.
      & test\_\allowbreak{}aufbereitung.py, Code~2, 9 \\
    \addlinespace
    Claude2026ar & Anthropic Claude Opus 5 & 13.08.2026
      & SHAP-Werte nach Lundberg und Lee zeigen nur, wie ein Modell seine Aufmerksamkeit verteilt. Ergänze deshalb eine Ablation je Faktorgruppe aus meiner Gruppierung, bei der das Modell ohne die Gruppe neu angepasst und bewertet wird.
      & m04\_\allowbreak{}shap.py \\
    \addlinespace
    Claude2026as & Anthropic Claude Opus 5 & 14.08.2026
      & Sorge dafür, dass ein zweiter Lauf dieselben Zahlen liefert: ein fester Startwert für alle Zufallsschritte, Bewertung auf einem Kern und eine Datei mit den genauen Paketversionen. Die Konstanten der Modellierung kommen in eine eigene Konfiguration.
      & config\_\allowbreak{}modelle.py, requirements\_\allowbreak{}lauf.txt \\
    \addlinespace
    Claude2026at & Anthropic Claude Opus 5 & 15.08.2026
      & Ob das Budget der Hyperparametersuche reicht, will ich wie von Bischl et~al. empfohlen prüfen. Wiederhole die Suche mit steigender Zahl an Versuchen, zeige den Verlauf als Kurve und ergänze einen kurzen Probelauf.
      & suchdiagnose.py \\
    \addlinespace
    Claude2026au & Anthropic Claude Opus 5 & 17.08.2026
      & Ein Macro-F1 lässt sich ohne Obergrenze schlecht einordnen. Schätze mit meinem angehängten Ansatz, wie gut die Einsatzart mit diesen Merkmalen höchstens vorhersagbar ist, und setze die Werte aus dem Vergleich ins Verhältnis dazu.
      & v4\_\allowbreak{}decke.py \\
    \addlinespace
    Claude2026av & Anthropic Claude Opus 5 & 20.08.2026
      & Zum Data Understanding nach Chapman et~al. brauche ich zwei getrennte Skripte: eine Merkmalstabelle mit Skalenniveau, Einheit und Quelle aus meiner Vorlage und eine deskriptive Auswertung mit Verteilung, Varianzzerlegung und Korrelationen.
      & codebook.py, deskriptiv.py \\
    \addlinespace
    Claude2026aw & Anthropic Claude Opus 5 & 23.08.2026
      & Der Census-Schlüssel steht im Klartext in der Konfiguration. Lies ihn stattdessen aus einer Umgebungsvariable.
      & config.py \\
    \addlinespace
    Claude2026ax & Anthropic Claude Opus 5 & 25.08.2026
      & Stelle aus den Rohquellen die Befunde zusammen, aus denen ein Eingriff in der Aufbereitung folgt, etwa Dubletten, fehlende Baujahre oder Lücken in der Abdeckung. Meine Liste der Eingriffe liegt bei, alles andere bleibt draußen.
      & rohbefunde.py \\
    \addlinespace
    Claude2026ay & Anthropic Claude Opus 5 & 29.08.2026
      & In den ACS-Jahrgängen mit älteren Zensusgrenzen fehlt ein Teil der Wohnbevölkerung. Suche die Ursache in der Zuordnung der Tracts, schlage eine Korrektur vor und lass die Aufbereitung abbrechen, wenn zu wenig Bevölkerung zugeordnet ist.
      & s1\_\allowbreak{}daten.py, Code~3 \\
    \addlinespace
    Claude2026az & Anthropic Claude Opus 5 & 30.08.2026
      & Bevor ich Ergebnisse ansehe, will ich wissen, welche Stadtteile im Hold-out und welche in der Entwicklung liegen. Erstelle ein Profil beider Hälften mit Klassenverteilung und Zielgrößen, unabhängig von jedem Modellergebnis.
      & panelprofil.py \\
    \addlinespace
    Claude2026ba & Anthropic Claude Opus 5 & 02.09.2026
      & Probst et~al. zeigen, dass Verfahren unterschiedlich stark auf ihre Hyperparameter reagieren. Setze jeden getunten Parametersatz auf allen anderen Folds ein und fasse zusammen, wie weit die Güte dabei schwankt.
      & parameter\allowbreak{}sensitivitaet.py \\
    \addlinespace
    Claude2026bb & Anthropic Claude Opus 5 & 11.09.2026
      & Nach Altman und Bland ist ein nicht signifikanter Test kein Beleg für Gleichheit. Berechne für die gepaarten Vergleiche Effektstärke und Trennschärfe bei zehn Wiederholungsmitteln, mit derselben Formel für die Streuung wie in den Vergleichsskripten.
      & trennschaerfe.py \\
    \addlinespace
    Claude2026bc & Anthropic Claude Opus 5 & 14.09.2026
      & Setze die README für die Abgabe auf, gegliedert nach Umgebung, Daten, Ablauf, Reproduzierbarkeit und Aufbau. Den Ablauf gebe ich als Befehlsfolge in einer Datei vor. Alle weiteren Inhalte stammen aus Code und Konfiguration, dazu ein Hinweis, dass ein neuer Abruf die Rohdaten überschreibt und der Census-Schlüssel nicht beiliegt.
      & README.md \\
    \addlinespace
    Claude2026bd & Anthropic Claude Opus 5.5 & 24.09.2026
      & Gleiche die Kommentare aller Skripte mit dem Stand der Arbeit ab. Meine angehängte Liste nennt veraltete Begriffe, Zahlen und Kapitelverweise. Schlage je Datei die Änderungen vor, geändert werden nur Kommentare und Docstrings.
      & alle Skripte \\
    \addlinespace
    Claude2026be & Anthropic Claude Opus 5.5 & 24.09.2026
      & Bringe die Kommentare in allen Dateien in eine einheitliche Form nach meiner Vorlage: ein Kopf mit Zweck, Aufruf, Eingang und Ausgang, darunter kurze Stichpunkte zu den Entscheidungen. Interne Kürzel und Verweise auf Arbeitsnotizen entfallen.
      & alle Skripte \\
    \addlinespace
\end{xltabular}
% \tabquelle
\endgroup
